\documentclass[aspectratio=169,9pt]{beamer}

% ------------------------------------------------------------------
% Theme: colors, fonts, helper macros, card styles, chrome
% (kept in its own package so it can be reused/edited independently)
% ------------------------------------------------------------------
\usepackage{style}
\usetikzlibrary{arrows.meta}

\newcommand{\R}{\mathbb{R}}

% Beamer's multimedia support: lets \movie embed a video that autostarts and
% loops when played in a PDF viewer that supports it (e.g. Pympress).
% Replaces the deprecated movie15/media9 route — no conflict with hyperref.
\usepackage{multimedia}

% ------------------------------------------------------------------
% Per-deck settings — adjust these here, not in style.sty, since they
% depend on where this particular deck's assets live and how you are
% currently using it (rehearsing with notes vs. exporting for the audience).
% ------------------------------------------------------------------
% Look for figures both in the current folder and common subfolders.
% Listed both relative to the repo root AND one level up (../...), so
% \includegraphics resolves whether latexmk is invoked from the repo root
% (the documented way — see .latexmkrc / the build command in the README-ish
% comments) or from inside slides/ (e.g. an editor that compiles "this file"
% from its own directory). NOTE: this only fixes *compilation*. For the
% embedded \movie below to actually find its video at presentation time, the
% final deck.pdf must end up at the repo root — see the comment by \movie.
\graphicspath{%
  {./}{../}%
  {assets/}{../assets/}%
  {pictures/}{../pictures/}%
  {figures/}{../figures/}%
  {figures/regenerated/}{../figures/regenerated/}%
  {figures/from-thesis/}{../figures/from-thesis/}%
}

% Let \input find files in figures/ directly (e.g. % Self-contained TikZ snippet — single-hidden-layer ReLU MLP on S^1.
% No preamble/document: meant to be pulled into a slide with \input.
% Uses the deck's palette (text, muted, line, border, tudcyan, bg) and
% \familydefault (sans), so it inherits the house style automatically.
\begin{tikzpicture}[
		x=1mm, y=1mm,
		every node/.style={font=\tiny},
		neuron/.style={circle, draw=text, fill=bg, minimum size=2mm, line width=0.5pt},
		outnode/.style={circle, draw=tudcyan, fill=bg, minimum size=2mm, line width=0.8pt},
		conn/.style={draw=text, line width=0.3pt},
		lbl/.style={font=\scriptsize, text=muted}
	]
	% --- node positions (mm); kept compact: natural size ~51mm x ~25mm,
	% i.e. roughly 38% of slide width and ~30mm tall (incl. labels) as required ---
	\coordinate (i1) at (0, 5);
	\coordinate (i2) at (0, -5);
	\coordinate (h1) at (14, 10);
	\coordinate (h2) at (14, 4);
	\coordinate (h3) at (14, -4);
	\coordinate (h4) at (14, -10);
	\coordinate (o)  at (28, 0);

	% --- edges first, so the nodes sit cleanly on top of them ---
	\foreach \a in {i1, i2}
	\foreach \b in {h1, h2, h3, h4}
	\draw[conn] (\a) -- (\b);
	\foreach \b in {h1, h2, h3, h4}
	\draw[conn] (\b) -- (o);

	% --- nodes (drawn over the edges; fill=bg masks the stubs cleanly) ---
	\node[neuron] at (i1) {};
	\node[neuron] at (i2) {};
	\node[neuron] at (h1) {};
	\node[neuron] at (h2) {};
	\node[text=text, font=\small] at (14, 1) {$\vdots$};
	\node[neuron] at (h3) {};
	\node[neuron] at (h4) {};
	\node[outnode] at (o) {};

	% --- labels (absolute coords, so no \usetikzlibrary{calc} is needed and
	% the snippet stays self-contained against the deck's default tikz setup) ---
	\node[lbl, anchor=east] at (-2, 0) {$x(\varphi)$};
	\node[lbl, anchor=west, text=tudcyan] at (30.5, 0) {$f(x)$};
	\node[lbl, align=center] at (14, -14) {$m$ neurons};
\end{tikzpicture}

\par\vspace{1mm}
{\scriptsize\itshape\textcolor{muted}{%
		$\displaystyle f_m(x;\theta)=\frac{1}{\sqrt m}\sum_{i=1}^m a_i\,\sigma\!\left(w_i^\top x + b_i\right),\ \ \sigma=\mathrm{ReLU}$%
	}}
),
% mirroring the dual-path \graphicspath above so it resolves whether latexmk
% runs from the repo root or from inside slides/.
\makeatletter
\def\input@path{{figures/}{../figures/}}
\makeatother

% Speaker notes: show them on a second screen while presenting.
% Swap to \setbeameroption{hide notes} for an audience-facing export.
% \setbeameroption{show notes on second screen=right}
\setbeameroption{hide notes}

% ------------------------------------------------------------------
\begin{document}

% ------------------------------------------------------------------
% Title slide
% ------------------------------------------------------------------
\begin{frame}[plain,noframenumbering]
	\vspace{3mm}
	\begin{columns}[T,onlytextwidth]
		\begin{column}{0.62\textwidth}
			\vspace{3em}
			{\fontsize{15}{27}\selectfont\bfseries Training Dynamics of\\Overparameterized Neural Networks}\par\vspace{2.5mm}
			\vspace{5mm}
			{\large Shreyas Kalvankar}\par\vspace{1mm}
			{\small MSc Computer Science, TU Delft}\par
			\vspace{6mm}
			{\scriptsize\textcolor{muted}{\textbf{Advisor:} Dr. D. M. J. Tax, Pattern Recognition \& Bioinformatics Group
					\\ \textbf{Supervisor:} Dr. A. Heinlein, Numerical Analysis Group \\ \textbf{Committee Member:} Prof. Dr. A. Papapantoleon, Applied Probability Group
				}}\par
		\end{column}
		\begin{column}{0.32\textwidth}
			\vspace{3mm}
			\begin{whitecard}[height=60mm]
				\centering
				\tudlogo[1.1cm]\par\vspace{7mm}
				\includegraphics[width=\linewidth]{cover_alt}
				\pill{Spectral Bias}{bluebg}{blue}
				\pill{NTK on $S^1$}{greenbg}{green}\par\vspace{0.2em}
				\pill{Fourier Analysis}{yellowbg}{orange}\par\vspace{0.2em}
				\pill{Finite-Scale Bounds}{purplebg}{purple}
			\end{whitecard}
		\end{column}
	\end{columns}
	\vfill
	{\color{line}\rule{\linewidth}{0.35pt}}\vspace{1mm}
	{\tiny\textcolor{muted}{May 2026}\hfill\textcolor{muted}{TU Delft}}
\end{frame}

% ------------------------------------------------------------------
% Learning means fitting a function to data (supervised-learning framework).
% Reveals in 4 overlay steps; the figure (figures/fit_intuition.tex) stages
% dots on step 1 and the fitted curve on step 2 via its overlay hooks.
%
% These define the figure's staging hooks: dots from step 1, curve from
% step 2. Defined OUTSIDE the frame (a raw #1 inside a frame body would be
% swallowed by beamer's argument scan) and BEFORE the \input, so the
% figure's own \providecommand identity-defaults are skipped. fit_intuition
% is used only on this slide, so leaving these defined is harmless.
\newcommand{\fitdotsoverlay}[1]{\visible<2->{#1}}
\newcommand{\fitcurveoverlay}[1]{\visible<4->{#1}}
\newcommand{\targetf}{{\color{orange}f^\star}}
\newcommand{\modelf}{{\color{tudcyan}f}}
\begin{frame}
	\note<1->{
		A neural network is just a rule that turns an input into an output,
		which is what we mean by a function. We are given example input output
		pairs and we want a function that agrees with them and, more
		importantly, gives sensible answers for inputs we have not seen, which
		is what generalize means. The way we get there is to adjust the
		function until its outputs on the examples are close to the true
		outputs, and the squared error written here just measures how far off
		it is. But we cannot try arbitrary functions one by one, so instead we
		fix a family of functions described by a set of adjustable numbers, the
		parameters, and learning means tuning those numbers. Keep the word
		unknowns in mind: the next motivation slide counts parameters against
		data points exactly like unknowns against equations. The next slide
		looks at what such a function is actually built from.
	}
	\vspace{0.5em}
	\slidetitle{INTRODUCTION}{Supervised Learning: Fitting data to a function}

	\vfill
	\begin{columns}[c,onlytextwidth]
		\begin{column}{0.52\textwidth}
			\setlength{\parskip}{2.2mm}\small

			\visible<2->{{\color{blue}\bfseries Data:}\ \ $n$ observations $(x_i, \targetf(x_i))$, examples of an unknown target function $\targetf$.\par}

			\visible<3->{{\color{green}\bfseries Model:}\ \ a learned prediction function $\modelf(x)$.\par}

			\visible<4->{{\color{blue}\bfseries Aim:}\ \ predict the target output $\targetf(x)$ for a new, unseen input $x$.\par}

			\visible<5->{{\color{orange}\bfseries How:}\ \ choose $\modelf$ so its predictions are close to the target function on the observed data, i.e.\ minimise
				$\mathcal{L}(f) = \displaystyle\frac{1}{n}\sum_{i=1}^{n}\big(\underbrace{\modelf(x_i)-\targetf(x_i)}_{:=\,r_i}\big)^2$.\par}

			\visible<6->{{\color{green}\bfseries In practice:}\ \ we cannot
				search over all possible functions, so we fix a parameterized
				family ${\color{tudcyan}f_\theta}(x)$ and learn by adjusting the
				parameters $\theta$. The parameters are the unknowns the data must
				pin down.\par}
		\end{column}
		\begin{column}{0.40\textwidth}
			\centering
			% Self-contained TikZ snippet — "examples become a function".
% No preamble/document: meant to be pulled into a slide with \input (bare
% filename, matching figures/mlp_diagram.tex; \input@path includes figures/).
% Abstract regression picture: scattered dots + one smooth curve threading
% near them. Uses the deck's palette (text, muted, tudcyan, line, border).
%
% STAGING HOOKS: the dots and the curve are each wrapped in a one-argument
% macro (\fitdotsoverlay / \fitcurveoverlay) so a slide can stage them with
% Beamer overlays, e.g.
%     \renewcommand{\fitdotsoverlay}[1]{\visible<1->{#1}}
%     \renewcommand{\fitcurveoverlay}[1]{\visible<2->{#1}}
%   BEFORE \input-ing this file. They default to identity (show everything),
% so the figure stays reusable as a plain static picture and no overlay specs
% are hardcoded here. The axes/frame are always drawn (visible from step 1).
\providecommand{\fitdotsoverlay}[1]{#1}
\providecommand{\fitcurveoverlay}[1]{#1}
\begin{tikzpicture}[
		x=1mm, y=1mm,
		dot/.style={circle, fill=text, inner sep=0pt, minimum size=1.5mm},
		axis/.style={draw=border, line width=0.5pt, ->}
	]
	% --- plain, unlabelled axes (always visible) ---
	\draw[axis] (0, 0) -- (44, 0);   % x-axis
	\draw[axis] (0, 0) -- (0, 30);   % y-axis
	\node[text=muted, font=\scriptsize, anchor=north] at (42, -2.2) {$x$};
	\node[text=muted, font=\scriptsize, anchor=east] at (-1.2, 29) {$y$};

	% --- the fitted curve (smooth, tudcyan) threading near the dots ---
	% Drawn before the dots so the dots sit on top when both are shown.
	\fitcurveoverlay{%
		\draw[tudcyan, line width=1pt, smooth, tension=0.7]
		plot coordinates {(3,7) (10,15) (18,21) (26,17) (33,11) (40,16)};
		\node[text=muted, font=\scriptsize, anchor=north] at (17, 26) {$f(x)$};
	}

	% --- ~9 scattered data points sitting NEAR the curve, not exactly on it ---
	\fitdotsoverlay{%
		\node[dot] at (4, 9)   {};
		\node[dot] at (9, 13)  {};
		\node[dot] at (13, 18) {};
		\node[dot] at (18, 19) {};
		\node[dot] at (22, 20) {};
		\node[dot] at (26, 15) {};
		\node[dot] at (31, 13) {};
		\node[dot] at (36, 12) {};
		\node[dot] at (41, 18) {};
	}
\end{tikzpicture}

		\end{column}
	\end{columns}
	\vfill
\end{frame}

% ------------------------------------------------------------------
% Motivation: overparameterization via the linear-equations analogy.
% Takeaway: many parameter settings fit the same data, so the interesting
% object is the choice gradient descent makes among them. The coarse-first
% teaser is confirmed by the demos on the next three slides.
\begin{frame}
	\note<1->{
		\begin{itemize}
			\item 3 eq / 3 unknowns: one solution; 3 eq / 4 unknowns: a line of solutions.
			\item Network: each data point = one constraint on $\theta$.
			\item ``Solution'' = near-zero training error.
			\item Example: CIFAR-10 has 50k images; WRN-28-10 has many more parameters.
			\item Large networks can generalize, but also fit random labels.
			\item Key question: which fit does gradient descent pick?
			\item Teaser: coarse before fine.
		\end{itemize}
	}
	\vspace{0.5em}
	\slidetitle{MOTIVATION \& RELEVANCE}{Overparameterization: many ways to fit the same data}
	{\textcolor{muted}{When many solutions fit the data, the learning process itself becomes important.}}

	\vspace{2mm}
	\begin{columns}[T,onlytextwidth]
		\begin{column}{0.43\textwidth}
			\begin{bluecard}[height=30mm]
				{\scriptsize\textcolor{blue}{\bfseries A linear analogy}}\par\vspace{2mm}
				\centering
				{\scriptsize
					\begin{tabular}{r@{\;\;$\Rightarrow$\;\;}l}
						$3$ equations, $3$ unknowns & one solution    \\[1.4mm]
						$3$ equations, $4$ unknowns & infinitely many \\
					\end{tabular}
				}\par\vspace{0.7mm}
				{\tiny\textcolor{muted}{independent equations, consistent system}}\par\vspace{2mm}
				{\scriptsize\raggedright
					One spare unknown, and the data no longer pins down the answer: a whole line of solutions fits.\par}
			\end{bluecard}
		\end{column}

		\begin{column}{0.53\textwidth}
			\visible<2->{
				\begin{orangecard}[height=30mm]
					{\scriptsize\textcolor{orange}{\bfseries The same situation in a neural network}}\par\vspace{1.2mm}

					{\scriptsize
						Each training point is one constraint on the parameters $\theta$:
						\[
							f_\theta(x_i)\approx y_i, \text{ for } i=1,\ldots,n.
						\]

						A \emph{solution}: a parameter setting with near-zero training error.\par\vspace{1.2mm}

						CIFAR-10 has \textbf{50k} training images
						(\textcolor{muted}{\href{https://www.cs.toronto.edu/~kriz/cifar.html}{Krizhevsky, 2009}}); \\
						WRN-28-10 has \textbf{36.5M} parameters and about \textbf{4\%} test error
						(\textcolor{muted}{\href{https://arxiv.org/abs/1605.07146}{Zagoruyko \& Komodakis, 2016}}).\par\vspace{1.4mm}
					}
				\end{orangecard}
			}
		\end{column}
	\end{columns}

	\vspace{1.5mm}
	\visible<3->{
		\begin{whitecard}
			\centering
			{\scriptsize
				Large networks can also fit random labels, so fitting the training set alone does not explain generalization
				(\textcolor{muted}{\href{https://arxiv.org/abs/1611.03530}{Zhang et al., 2017}}).\\
				Empirical observation: models often learn simple structure before complex detail
				(\textcolor{muted}{\href{https://proceedings.mlr.press/v97/rahaman19a.html}{Rahaman et al., 2019}};
				\textcolor{muted}{\href{https://arxiv.org/abs/2410.02348}{Boursier \& Flammarion, 2024}}).
			}
		\end{whitecard}

		% \vspace{0.3mm}
		% {\tiny\textcolor{muted}{
		% 		\href{https://www.cs.toronto.edu/~kriz/cifar.html}{\texttt{cs.toronto.edu/\textasciitilde kriz/cifar.html}}
		% 		\quad
		% 		\href{https://arxiv.org/abs/1605.07146}{\texttt{arxiv.org/abs/1605.07146}}
		% 	}}\par\vspace{0.1mm}
		% {\tiny\textcolor{muted}{
		% 		\href{https://arxiv.org/abs/1611.03530}{\texttt{arxiv.org/abs/1611.03530}}
		% 		\quad
		% 		\href{https://proceedings.mlr.press/v97/rahaman19a.html}{\texttt{proceedings.mlr.press/v97/rahaman19a.html}}
		% 	}}
	}
\end{frame}

% ------------------------------------------------------------------
% Broader relevance: not reliability, not a deployment claim.
% Takeaway: understanding training behaviour is a foundation for making
% overparameterized models less trial-and-error driven.
\begin{frame}
	\note<1->{
		\begin{itemize}
			\item Keep this high level.
			\item Not claiming this solves energy or deployment.
			\item Point: training behaviour matters because training is costly and opaque.
			\item Compute and data-centre energy are growing quickly.
			\item Do not explain thesis yet; next slides show the mechanism.
			\item Transition: what does training learn first?
		\end{itemize}
	}
	\vspace{0.5em}
	\slidetitle{MOTIVATION \& RELEVANCE}{Why study training behaviour?}
	{\textcolor{muted}{The broader goal is to understand learning, not just report a final test error.}}

	\vspace{3mm}
	{\small
		\begin{itemize}
			\setlength{\itemsep}{3.5mm}

			\item Once many fits are possible, improving a model is not only about making it larger; it is also about understanding what training is actually doing.

			\item This matters more as modern AI systems become expensive and energy-intensive to train:
			      the compute used to train large-scale AI models is reported to double roughly every five months
			      (\textcolor{muted}{\href{https://hai.stanford.edu/ai-index/2025-ai-index-report/research-and-development}{Stanford AI Index, 2025}}); the Llama~3.1 family alone used 39.3M H100 GPU-hours for training (\textcolor{muted}{\href{https://github.com/meta-llama/llama-models/blob/main/models/llama3_1/MODEL_CARD.md}{Meta AI, 2024}}).

			\item So before asking for a bigger model, we can ask: \textbf{what does gradient descent learn first?}
		\end{itemize}
	}
\end{frame}

% ------------------------------------------------------------------
\begin{frame}
	\note<1->{
		\begin{itemize}
			\item Target = short sum of frequencies, coarse $\to$ fine.
			\item The object the rest of the talk tracks.
			\item Sets up the demo: what appears first?
		\end{itemize}
	}
	\vspace{0.5em}
	\slidetitle{ILLUSTRATION}{A concrete target}

	\centering
	{\textcolor{muted}{To watch that choice happen, take a target built from a coarse shape plus finer and finer detail.}}

	\vspace{1.5mm}
	\begin{columns}[T,onlytextwidth]
		\begin{column}{0.61\textwidth}
			\centering
			% The path below is relative to the FINAL deck.pdf's location, not
			% to this .tex file or the compiling directory — Pympress resolves
			% it relative to the PDF itself. It is only correct if deck.pdf
			% ends up at the repo root (i.e. compiled with
			% `latexmk ... slides/deck.tex` run FROM the repo root, per
			% .latexmkrc — the documented way). If deck.pdf instead lands in
			% slides/, the movie will not be found at presentation time even
			% though the slide compiles and looks fine.
			\movie[autostart, loop, poster, showcontrols=false]{%
				\includegraphics[width=\linewidth]{target_construction_poster.png}%
			}{figures/regenerated/target_construction.mp4}
		\end{column}
		\begin{column}{0.37\textwidth}
			\begin{whitecard}
				{\bfseries Frequency components}\par\vspace{1.5mm}
				\begin{itemize}\tightlist\scriptsize
					\item Each piece (i.e. a \emph{mode}) is one frequency, from slow waves to fast ones.
					\item Low frequencies set the coarse shape; high frequencies add the fine detail.
					\item The target is a short sum of these pieces.
				\end{itemize}
			\end{whitecard}
		\end{column}
	\end{columns}

	\pause
	\vspace{1.5mm}
	\begin{bluecard}[width=0.6\linewidth]
		\centering{\Large\bfseries What gets learned first?}
	\end{bluecard}
\end{frame}

% ------------------------------------------------------------------
% Slide A — the phenomenon (single-hidden-layer MLP, fixed split)
% ------------------------------------------------------------------
\begin{frame}
	\note<1->{
		We train the single-hidden-layer network, the model the whole thesis
		analyses, on the target from the opener. The three panels show the fit
		climbing toward the target, how each mode's residual decays during
		training, and the output rebuilt as a partial sum, coarse to fine. The
		takeaway: low-frequency structure is learned before high-frequency
		detail.
	}
	\vspace{0.5em}
	\slidetitle{ILLUSTRATION}{Watch one network learn it}
	{\textcolor{muted}{Train a single-hidden-layer network on that target and watch the residual: what gets learned first?}}

	\vspace{1.5mm}
	\begin{columns}[T,onlytextwidth]
		\begin{column}{0.59\textwidth}
			\centering
			\movie[autostart, loop, poster, showcontrols=false]{%
				\includegraphics[width=\linewidth]{spectral_bias_mlp_single_layer_slow_intro_poster.png}%
			}{figures/regenerated/spectral_bias_mlp_single_layer_slow_intro.mp4}
		\end{column}
		\begin{column}{0.39\textwidth}
			\centering
			% Self-contained TikZ snippet — single-hidden-layer ReLU MLP on S^1.
% No preamble/document: meant to be pulled into a slide with \input.
% Uses the deck's palette (text, muted, line, border, tudcyan, bg) and
% \familydefault (sans), so it inherits the house style automatically.
\begin{tikzpicture}[
		x=1mm, y=1mm,
		every node/.style={font=\tiny},
		neuron/.style={circle, draw=text, fill=bg, minimum size=2mm, line width=0.5pt},
		outnode/.style={circle, draw=tudcyan, fill=bg, minimum size=2mm, line width=0.8pt},
		conn/.style={draw=text, line width=0.3pt},
		lbl/.style={font=\scriptsize, text=muted}
	]
	% --- node positions (mm); kept compact: natural size ~51mm x ~25mm,
	% i.e. roughly 38% of slide width and ~30mm tall (incl. labels) as required ---
	\coordinate (i1) at (0, 5);
	\coordinate (i2) at (0, -5);
	\coordinate (h1) at (14, 10);
	\coordinate (h2) at (14, 4);
	\coordinate (h3) at (14, -4);
	\coordinate (h4) at (14, -10);
	\coordinate (o)  at (28, 0);

	% --- edges first, so the nodes sit cleanly on top of them ---
	\foreach \a in {i1, i2}
	\foreach \b in {h1, h2, h3, h4}
	\draw[conn] (\a) -- (\b);
	\foreach \b in {h1, h2, h3, h4}
	\draw[conn] (\b) -- (o);

	% --- nodes (drawn over the edges; fill=bg masks the stubs cleanly) ---
	\node[neuron] at (i1) {};
	\node[neuron] at (i2) {};
	\node[neuron] at (h1) {};
	\node[neuron] at (h2) {};
	\node[text=text, font=\small] at (14, 1) {$\vdots$};
	\node[neuron] at (h3) {};
	\node[neuron] at (h4) {};
	\node[outnode] at (o) {};

	% --- labels (absolute coords, so no \usetikzlibrary{calc} is needed and
	% the snippet stays self-contained against the deck's default tikz setup) ---
	\node[lbl, anchor=east] at (-2, 0) {$x(\varphi)$};
	\node[lbl, anchor=west, text=tudcyan] at (30.5, 0) {$f(x)$};
	\node[lbl, align=center] at (14, -14) {$m$ neurons};
\end{tikzpicture}

\par\vspace{1mm}
{\scriptsize\itshape\textcolor{muted}{%
		$\displaystyle f_m(x;\theta)=\frac{1}{\sqrt m}\sum_{i=1}^m a_i\,\sigma\!\left(w_i^\top x + b_i\right),\ \ \sigma=\mathrm{ReLU}$%
	}}


			\vspace{2mm}
			\begin{whitecard}
				{\bfseries Notice}\par\vspace{1mm}
				\begin{itemize}\tightlist\scriptsize
					\item The fit climbs toward the target.
					\item Low-frequency modes die first.
					\item Fine detail fills in last.
				\end{itemize}
			\end{whitecard}
		\end{column}
	\end{columns}

	\vspace{-8mm}
	\begin{bluecard}[width=0.55\linewidth]
		\centering{\bfseries The coarse shape appears first; the fine detail comes last.}
	\end{bluecard}
\end{frame}

% ------------------------------------------------------------------
\begin{frame}
	\note<1->{
		Spectral bias is the tendency of neural networks trained with gradient descent to learn low-frequency structure before high-frequency detail.

		In other words, the network first captures the smooth, global shape of the target, and only later fits the more oscillatory parts.

		That is what these examples are meant to show. Even though the architectures are different, the learning process has the same qualitative pattern: low-frequency modes appear early, high-frequency modes are delayed.

		My thesis asks whether we can explain this from the training dynamics.

		The main idea is to reformulate spectral bias as an operator eigenvalue problem. If an operator governs the residual dynamics, then its eigenvalues determine which frequencies are learned quickly and which ones are learned slowly.

		So the three guiding questions are: what controls mode-wise learning rates, does the explanation survive finite data, and does it survive finite width?
	}
	\vspace{0.5em}
	\slidetitle{THESIS}{Why does gradient descent learn simple structure first?}
	{\small\textcolor{muted}{This pattern has a name: the frequency principle, also called spectral bias (\href{https://arxiv.org/abs/1901.06523}{Xu et al., 2019}).}}

	\vspace{1.8mm}
	\begin{columns}[T,onlytextwidth]
		\begin{column}{0.48\textwidth}
			\begin{bluecard}[height=22mm]
				\textcolor{purple}{\scriptsize\bfseries Thesis in one sentence}\par\vspace{1mm}
				{\scriptsize Derive the exact learning speed of every frequency
					in an idealized network, then quantify how far that exact
					picture persists in a realistic one: finitely many training
					points, finitely many neurons.}
			\end{bluecard}
		\end{column}
		\begin{column}{0.50\textwidth}
			\begin{whitecard}[height=23mm]
				{\scriptsize\bfseries Three questions}\par\vspace{1.3mm}
				\questionnum{blue}{1}\quad {\tiny What sets the learning speed of each frequency?}\par\vspace{1.2mm}
				\questionnum{green}{2}\quad {\tiny Does the answer survive when the network sees only \(n\) sampled points? \emph{(RQ1)}}\par\vspace{1.2mm}
				\questionnum{orange}{3}\quad {\tiny Does it survive with finitely many neurons, at the start of training and during it? \emph{(RQ2)}}
			\end{whitecard}
		\end{column}
	\end{columns}

	\vspace{2mm}
	\begin{whitecard}
		\begin{minipage}{\linewidth}
			\raisebox{1.4mm}{\scriptsize\bfseries Frequency ordering in different networks.}
			\hfill
			\includegraphics[width=0.40\linewidth]{legend.png}
		\end{minipage}
		\begin{columns}[T,onlytextwidth]
			\begin{column}{0.32\textwidth}
				\centering
				\movie[autostart, loop, poster, showcontrols=false]{%
					\includegraphics[width=0.86\linewidth]{spectral_bias_mlp_single_layer_poster.png}%
				}{figures/regenerated/spectral_bias_mlp_single_layer.mp4}\par\vspace{0.5mm}
				{\tiny One hidden layer}
			\end{column}

			\begin{column}{0.32\textwidth}
				\centering
				\movie[autostart, loop, poster, showcontrols=false]{%
					\includegraphics[width=0.86\linewidth]{spectral_bias_mlp_deep_2_layer_poster.png}%
				}{figures/regenerated/spectral_bias_mlp_deep_2_layer.mp4}\par\vspace{0.5mm}
				{\tiny Two hidden layers}
			\end{column}

			\begin{column}{0.32\textwidth}
				\centering
				\movie[autostart, loop, poster, showcontrols=false]{%
					\includegraphics[width=0.86\linewidth]{spectral_bias_attention_block_poster.png}%
				}{figures/regenerated/spectral_bias_attention_block.mp4}\par\vspace{0.5mm}
				{\tiny Attention Block}
			\end{column}
		\end{columns}
	\end{whitecard}
\end{frame}

\begin{frame}
	\note<1->{
		Picture the error as a landscape: high where the network is wrong, low
		where it fits. Training rolls downhill, taking the steepest way down, which
		is why it is called gradient descent. Notice the steps shrink as it nears
		the bottom: the closer the error gets to zero, the smaller the correction,
		because the slope itself is proportional to how much error is left. Then we
		look at the same bowl from straight above. The rings are contours, lines of
		equal height, so closely spaced rings mean steep and widely spaced rings
		mean shallow. This bowl has a steep direction and a shallow one, and that is
		the distinction the next slide builds on.
	}
	\vspace{0.5em}
	\slidetitle{BACKGROUND}{The error is a landscape}
	{\textcolor{muted}{Seen from above, the contour rings are the bowl's height; training rolls toward the bottom.}}

	{\raggedleft\tiny\textcolor{muted}{recall the loss: $\mathcal{L}(f) = \tfrac1n\sum_i\big(f(x_i)-f^\star(x_i)\big)^2$}\par}

	\vspace{2mm}
	\centering
	\movie[autostart, loop, poster, showcontrols=false]{%
		\includegraphics[width=0.62\linewidth]{bowl_to_contours_poster.png}%
	}{figures/regenerated/bowl_to_contours.mp4}
\end{frame}

% ------------------------------------------------------------------
% Keystone slide: "rate = eigenvalue", introduced geometrically before any
% NTK machinery. Self-contained TikZ (no external assets).
\begin{frame}
	\note<1->{
		One sentence on what training does: it makes the error small. It does this by
		walking downhill on the error, stepping in the direction the error drops
		fastest, which is the negative gradient, hence gradient descent. Now the
		question this slide answers: why do different directions get learned at
		different speeds. Picture the error landscape as a stretched bowl. In the
		bowl's own tilted coordinates the descent path mixes the two directions
		together, so you cannot read off a clean rate for either one. The fix is to
		use the bowl's own axes as coordinates. These special directions, the
		eigenvectors, are the ones the dynamics never mixes. Along a single one of
		them the error falls in proportion to how much is left, the differential
		equation here, which is the same law as a cooling cup of coffee or
		radioactive half-life, so it decays exponentially, alpha of t equals alpha of
		zero times e to the minus lambda t. The rate lambda is the eigenvalue, large
		is fast, small is slow. Everything after this works out what these directions
		and rates are for a real network on the circle.
	}
	\vspace{0.4em}
	\slidetitle{BACKGROUND}{Each direction has its own rate}
	{\textcolor{muted}{Training makes the error small by walking downhill, along the steepest drop.}}\par
	{\raggedleft\tiny\textcolor{muted}{recall the loss: $\mathcal{L}(f) = \tfrac1n\sum_i\big(f(x_i)-f^\star(x_i)\big)^2$}\par}

	\vspace{1mm}
	% ---------- ROW 1: the bowl, coupled then rotated to its own axes ----------
	\begin{columns}[c,onlytextwidth]
		\begin{column}{0.45\textwidth}
			\centering
			\visible<1->{
				\begin{tikzpicture}[scale=0.50]
					% Handed x, y axes (kept); the bowl is tilted RELATIVE to them.
					% Matches the lead-in "error is a landscape" frozen frame.
					\draw[line] (-2.5,0)--(2.5,0);
					\draw[line] (0,-1.7)--(0,1.7);
					\node[muted,font=\tiny] at (2.66,0.0) {$x$};
					\node[muted,font=\tiny] at (0.0,1.9) {$y$};
					% Tilted contour ellipses: principal axes at 30 deg, NOT along x, y.
					\foreach \r in {1.0,0.861,0.707,0.544,0.36}{
							\draw[tudcyan,opacity=0.55,rotate=30] (0,0) ellipse ({1.95*\r} and {1.36*\r});
						}
					% Curved gradient-descent path in x, y: each step crosses the
					% tilted contours, so the route bends (the coordinates are
					% coupled). Matches the lead-in path; starts high off the x axis,
					% the steep part dies first, then it slides along the shallow axis.
					\draw[orange,very thick]
					(0.792,1.358)--(0.679,0.734)--(0.523,0.432)--(0.384,0.271)--(0.276,0.178)--(0.196,0.120)--(0.138,0.082);
					\foreach \p in {(0.679,0.734),(0.523,0.432),(0.384,0.271),(0.276,0.178),(0.196,0.120),(0.138,0.082)}{
							\fill[orange] \p circle (1.2pt);
						}
					\fill[orange] (0.792,1.358) circle (1.7pt);
				\end{tikzpicture}\par\vspace{0.5mm}
				{\scriptsize\textcolor{muted}{In x, y axes the two directions are coupled.}}
			}
		\end{column}

		\begin{column}{0.08\textwidth}
			\centering
			\visible<2->{$\Longrightarrow$\\[0.5mm]{\tiny\textcolor{tudcyan}{the bowl's own axes}}}
		\end{column}

		\begin{column}{0.45\textwidth}
			\centering
			\visible<2->{
				\begin{tikzpicture}[scale=0.50]
					% The SAME bowl, rotated into its own (eigen) axes: the tilted
					% ellipse of panel 1 is now axis-aligned, so the directions
					% decouple. Same ellipse aspect (1.95 : 1.36) as panel 1.
					\draw[line] (-2.4,0)--(2.4,0);
					\draw[line] (0,-1.6)--(0,1.6);
					\foreach \r in {1.0,0.861,0.707,0.544,0.36}{
							\draw[tudcyan,opacity=0.55] (0,0) ellipse ({1.95*\r} and {1.36*\r});
						}
					% The SAME route in these (eigen) coordinates: the fast
					% (vertical) component dies first, then it slides in along the
					% slow (horizontal) axis. This is the decoupling made visible.
					\draw[orange,very thick]
					(1.365,0.780)--(0.956,0.296)--(0.669,0.113)--(0.468,0.043)--(0.328,0.016)--(0.229,0.006)--(0.161,0.002);
					\foreach \p in {(0.956,0.296),(0.669,0.113),(0.468,0.043),(0.328,0.016),(0.229,0.006),(0.161,0.002)}{
							\fill[orange] \p circle (1.1pt);
						}
					\fill[orange] (1.365,0.780) circle (1.6pt);
					\draw[blue!60,very thick,-{Latex}] (0,0)--(1.95,0);
					\node[blue!60!black,font=\tiny,anchor=north] at (1.4,-0.12) {slow: small $\lambda$};
					\draw[orange!80,very thick,-{Latex}] (0,0)--(0,1.36);
					\node[orange!80!black,font=\tiny,anchor=south west] at (0.06,1.36) {fast: large $\lambda$};
				\end{tikzpicture}\par\vspace{0.5mm}
				{\scriptsize\textcolor{muted}{Along each axis the error falls on its own.}}
			}
		\end{column}
	\end{columns}

	\vspace{1mm}
	% ---------- ROW 2: one decoupled direction, its ODE and exponential ----------
	\begin{columns}[T,onlytextwidth]
		\begin{column}{0.49\textwidth}
			\visible<3->{
				\begin{bluecard}[height=22mm]
					{\scriptsize Along one direction, the error falls in proportion to how much is left:}\par\vspace{1mm}
					\scriptsize\[ \frac{d}{dt}\alpha(t) = -\lambda\,\alpha(t) \]
				\end{bluecard}
			}
		\end{column}
		\begin{column}{0.49\textwidth}
			\visible<4->{
				\begin{greencard}[height=22mm]
					\begin{columns}[c,onlytextwidth]
						\begin{column}{0.58\textwidth}
							{\scriptsize so it decays exponentially:}\par\vspace{1mm}
							\[ \alpha(t)=\alpha(0)\,e^{-\lambda t} \]
							{\tiny like cooling or half-life; $\lambda$ is the rate.}
						\end{column}
						\begin{column}{0.40\textwidth}
							\centering
							\begin{tikzpicture}[scale=0.42]
								\draw[line] (0,0)--(3.2,0);
								\draw[line] (0,0)--(0,2.4);
								\node[muted,font=\tiny,anchor=north] at (3.05,-0.08) {$t$};
								\node[muted,font=\tiny,anchor=west] at (0.30,2.22) {$\alpha(t)$};
								\draw[green!60!black,very thick] plot[smooth] coordinates
									{(0,2.2)(0.5,1.47)(1.0,0.99)(1.5,0.66)(2.0,0.44)(2.6,0.27)(3.1,0.16)};
							\end{tikzpicture}
						\end{column}
					\end{columns}
				\end{greencard}
			}
		\end{column}
	\end{columns}

	\vspace{1mm}
	\centering
	\visible<5->{
		\begin{bluecard}[width=0.9\linewidth]
			\centering
			{\scriptsize {\bfseries Eigenvector:} a direction the dynamics keeps separate, never mixing it with others.\quad
				{\bfseries Eigenvalue} $\lambda$: its rate, large is fast, small is slow.}\par\vspace{0.8mm}
			{\bfseries What are these directions and rates for a neural network?}
		\end{bluecard}
	}
\end{frame}

% ------------------------------------------------------------------
% Bridge 1: the scalar rate becomes an operator (the NTK). Answers the keystone's
% closing question in principle; stays model-agnostic (no f_m yet, that is SETUP).
\begin{frame}
	\note<1->{
		Picking up the question from the last slide. For a single direction the error
		obeyed alpha dot equals minus lambda alpha, the rate was one number. For a
		network the error is a whole function, the residual, and it still falls in
		proportion to how much is left, but the proportion is no longer a number. It
		is an operator acting on the whole function, and that operator is the neural
		tangent kernel. The residual obeys r dot equals minus T r. This holds for any
		differentiable model under gradient flow, which is why I am not writing the
		network down yet. If T were a fixed operator, diagonalizing it would split
		this back into one copy of the scalar law per eigenfunction, alpha k dot
		equals minus lambda k alpha k. So the eigenfunctions are the directions and
		the eigenvalues are the rates, and the question from the last slide becomes:
		what is the spectrum of T.
	}
	\vspace{0.4em}
	\slidetitle{BACKGROUND}{The rate becomes an operator: the Neural Tangent Kernel (\textcolor{muted}{\href{https://arxiv.org/abs/1806.07572}{Jacot et al., 2018}})}
	{\textcolor{muted}{The same law as before, with the single rate replaced by an operator.}}

	{\raggedleft\tiny\textcolor{muted}{recall the residual: $r(x) = f(x)-f^\star(x)$}\par}

	\vspace{2mm}
	\begin{columns}[T,onlytextwidth]
		\begin{column}{0.49\textwidth}
			\begin{bluecard}[height=26mm]
				{\textcolor{blue}{\bfseries One direction}}\par\vspace{1.5mm}
				\[ \frac{d}{dt}\alpha(t) = -\lambda\,\alpha(t) \]
				{\scriptsize The error falls at rate $\lambda$, a single number.}
			\end{bluecard}
		\end{column}
		\begin{column}{0.49\textwidth}
			\begin{greencard}[height=26mm]
				{\textcolor{green}{\bfseries A whole function}}\par\vspace{1.5mm}
				\[ \partial_t r_t = -T_t\,r_t \]
				{\scriptsize The rate becomes an operator $T_t$, the network's tangent kernel (NTK).}
			\end{greencard}
		\end{column}
	\end{columns}

	\vspace{1.5mm}
	\begin{whitecard}
		\begin{columns}[c,onlytextwidth]
			\begin{column}{0.60\textwidth}
				{\scriptsize\raggedright If $T_t$ were time-independent, diagonalizing it
					would split this back into one copy of the scalar law per
					direction \(k\):\par}
			\end{column}
			\begin{column}{0.39\textwidth}
				\vspace{-1.5mm}
				\[ \dot\alpha_k = -\lambda_k\,\alpha_k \]
				\vspace{-4mm}
			\end{column}
		\end{columns}
		\vspace{1.5mm}
		{\centering\scriptsize
		{\bfseries eigenfunctions} are the directions
		\;\textcolor{muted}{$\cdot$}\;
		{\bfseries eigenvalues} $\lambda_k$ are the rates
		\;\textcolor{muted}{$\cdot$}\;
		together they form the {\bfseries spectrum} of the time-independent $T$\par}
	\end{whitecard}

	\centering
	\vspace{1mm}
	\begin{bluecard}[width=0.6\linewidth]
		\centering{\bfseries So the question becomes: what is the spectrum of $T$?}
	\end{bluecard}
\end{frame}

% ------------------------------------------------------------------
% Bridge 2: why the thesis idealizes before analyzing. Two obstructions + branching roadmap.
\begin{frame}
	\note<1->{
		To read off a fixed set of directions and rates, T has to be one fixed
		operator, and it is not yet. Two things are in the way. First, it is time
		dependent: T changes as the network trains, so its directions and rates
		drift. We freeze it to get a constant operator. Second, it depends on the
		scale: at finite data it depends on the sampled points, and at finite width
		it is random because it depends on the initialization. We take the continuum
		and infinite width, where it becomes a single deterministic operator with
		exact structure. The strategy for the rest of the talk: analyze this clean
		reference first, where the spectrum is exact, then relax each idealization
		separately. Finite data is one perturbation of the reference; finite width
		is the other, first with the kernel frozen, then letting it evolve.
	}
	\vspace{0.4em}
	\slidetitle{SETUP}{Making the dynamics analyzable}
	{\textcolor{muted}{To read off a fixed set of directions and rates, $T$ must be one fixed operator. Two things are in the way.}}

	{\raggedleft\tiny\textcolor{muted}{recall: $ \partial_t r_t = -T_t\,r_t $}\par}

	\vspace{2mm}
	\begin{columns}[T,onlytextwidth]
		\begin{column}{0.49\textwidth}
			\begin{orangecard}[height=28mm]
				{\textcolor{orange}{\bfseries Time-dependent}}\par\vspace{1.5mm}
				{\scriptsize $T_t$ changes as the network trains, so its directions and rates drift.}\par\vspace{1.5mm}
				{\scriptsize \textbf{Solution:} Take width to infinity to get a constant operator $T$
					(\textcolor{muted}{\href{https://arxiv.org/abs/1806.07572}{Jacot et al., 2018}}).}
			\end{orangecard}
		\end{column}
		\begin{column}{0.49\textwidth}
			\begin{purplecard}[height=28mm]
				{\textcolor{purple}{\bfseries Finite scale}}\par\vspace{1.5mm}
				{\scriptsize At finite data $T$ depends on the sampled input points, so it changes on what it sees.}\par\vspace{1.5mm}
				{\scriptsize \textbf{Solution:} Consider using the continuum as input.
					(\textcolor{muted}{\href{https://arxiv.org/abs/1806.07572}{Jacot et al., 2018}, \href{https://arxiv.org/abs/1912.01198}{Cao et al., 2019}}).}
			\end{purplecard}
		\end{column}
	\end{columns}

	\vspace{1mm}
	\begin{whitecard}
		\centering
		{\scriptsize\bfseries Strategy:}\ {\scriptsize build the clean reference, then relax each idealization separately.}\par\vspace{1.5mm}
		\begin{tikzpicture}[
				pillnode/.style={
						rounded corners=8pt,
						inner xsep=6pt,
						inner ysep=3pt,
						font=\scriptsize\bfseries
					}
			]
			\node[pillnode, fill=bluebg, text=blue] (ref) at (0,0) {Reference model};
			\node[pillnode, fill=yellowbg, text=orange] (data) at (3.6,0.45) {Finite data};
			\node[pillnode, fill=purplebg, text=purple] (width) at (3.6,-0.45) {Finite width};
			\node[pillnode, fill=greenbg, text=green] (evol) at (6.6,-0.45) {Evolving};

			\draw[muted,->] (ref.east) -- (data.west);
			\draw[muted,->] (ref.east) -- (width.west);
			\draw[muted,->] (width.east) -- (evol.west);
		\end{tikzpicture}
	\end{whitecard}
\end{frame}

% % ------------------------------------------------------------------
% % Bridge 2: why the thesis idealizes before analyzing. Two obstructions + branching roadmap.
% \begin{frame}
% 	\note<1->{
% 		To read off a fixed set of directions and rates, T has to be one fixed
% 		operator, and it is not yet. Two things are in the way, and they live in
% 		two different objects. The kernel Theta, defined up here as the inner
% 		product of parameter gradients, is an average over the m random neurons,
% 		so at finite width it is random, and it moves as the network trains. At
% 		infinite width the average concentrates at its expectation, and at this
% 		scaling the parameters barely move, so the kernel is deterministic and
% 		frozen. The operator T applies that kernel as an average over the n data
% 		points; in the continuum that sum becomes an integral over the circle.
% 		Notice both idealizations are the same move: replace an empirical average
% 		by its population version, over neurons in one card, over data points in
% 		the other. The finite-scale errors later in the talk are exactly the cost
% 		of undoing this, which is why both shrink like one over the square root
% 		of the scale. Strategy for the rest of the talk: analyze the clean
% 		reference, then relax each idealization separately; finite data is one
% 		perturbation, finite width the other, first frozen, then evolving.
% 	}
% 	\vspace{0.4em}
% 	\slidetitle{BACKGROUND}{Making the dynamics analyzable}
% 	{\textcolor{muted}{To read off a fixed set of directions and rates, $T$ must be one fixed operator. Two things are in the way.}}
%
% 	{\raggedleft\tiny\textcolor{muted}{recall: $\partial_t r_t = -T_t\,r_t$,\ $\Theta_t(x,x') = \big\langle \nabla_\theta f(x;\theta_t), \nabla_\theta f(x';\theta_t)\big\rangle$}\par}
%
% 	\vspace{1mm}
% 	\begin{columns}[T,onlytextwidth]
% 		\begin{column}{0.49\textwidth}
% 			\begin{orangecard}[height=32mm]
% 				{\textcolor{orange}{\bfseries The kernel $\Theta$: random and moving at finite width}}\par\vspace{1mm}
% 				{\scriptsize For our network, $\Theta_t$ is an average of one term per neuron:}\par\vspace{-2mm}
% 				{\tiny
% 					\[
% 						\Theta_t(x,x') = \frac{1}{m}\sum_{i=1}^{m}\kappa_i(x,x')
% 						\;\xrightarrow{\;m\to\infty\;}\;
% 						\mathbb{E}\big[\kappa(x,x')\big] = \Theta_\infty(x,x')
% 					\]
% 				}
% 				\vspace{-2.5mm}\par
% 				{\tiny At initialization the average concentrates at its expectation; at
% 					this scaling the parameters barely move during training, so the kernel
% 					stays frozen there\
% 					(\textcolor{muted}{\href{https://arxiv.org/abs/1806.07572}{Jacot et al., 2018}}).}
% 			\end{orangecard}
% 		\end{column}
% 		\begin{column}{0.49\textwidth}
% 			\begin{purplecard}[height=32mm]
% 				{\textcolor{purple}{\bfseries The operator $T$: built from finitely many samples}}\par\vspace{1mm}
% 				{\scriptsize At $n$ samples, $T$ averages the kernel over the data points:}\par\vspace{-2mm}
% 				{\tiny
% 					\[
% 						\frac{1}{n}\sum_{i=1}^{n}\Theta(x,\varphi_i)\,r(\varphi_i)
% 						\;\xrightarrow{\;n\to\infty\;}\;
% 						\int_{S^1}\Theta(x,\varphi)\,r(\varphi)\,d\mu(\varphi)
% 					\]
% 				}
% 				\vspace{-2.5mm}\par
% 				{\tiny The continuum replaces the sum over samples by an integral over
% 					the whole circle\
% 					(\textcolor{muted}{\href{https://arxiv.org/abs/1806.07572}{Jacot et al., 2018}};
% 					\textcolor{muted}{\href{https://arxiv.org/abs/1912.01198}{Cao et al., 2019}}).}
% 			\end{purplecard}
% 		\end{column}
% 	\end{columns}
% 	\begin{whitecard}
% 		\centering
% 		{\scriptsize\bfseries Strategy:}\ {\scriptsize build the clean reference, then relax each idealization separately.}\par\vspace{1.5mm}
% 		\begin{tikzpicture}
% 			\node (ref)   at (0,0)       {\pill{Reference model}{bluebg}{blue}};
% 			\node (data)  at (3.6,0.45)  {\pill{Finite data}{yellowbg}{orange}};
% 			\node (width) at (3.6,-0.45) {\pill{Finite width}{purplebg}{purple}};
% 			\node (evol)  at (6.6,-0.45) {\pill{Evolving}{greenbg}{green}};
% 			\draw[muted,-{Latex}] (ref.east) -- (data.west);
% 			\draw[muted,-{Latex}] (ref.east) -- (width.west);
% 			\draw[muted,-{Latex}] (width.east) -- (evol.west);
% 		\end{tikzpicture}
% 	\end{whitecard}
% \end{frame}


% ------------------------------------------------------------------
% Setup: the concrete model and domain. Keep this light because the NTK
% dynamics were introduced on the previous slides.
\begin{frame}
	\note<1->{
		\begin{itemize}
			\item Instantiate the abstract operator story.
			\item Circle gives Fourier modes: coarse to fine.
			\item Model: one-hidden-layer ReLU network.
			\item Define \(\phi_k\) as a Fourier mode and track coefficients \(\alpha_k(t)\).
			\item Next: ideal limit makes these exact eigen-directions.
		\end{itemize}
	}
	\vspace{0em}
	\slidetitle{SETUP}{The concrete model: regression on $S^1$}
	{\textcolor{muted}{The unit circle gives a natural language for ``coarse'' and ``fine'': Fourier modes.}}

	\vspace{0.2mm}
	\begin{columns}[T,onlytextwidth]
		\begin{column}{0.52\textwidth}
			{\scriptsize
				\begin{itemize}
					\setlength{\itemsep}{1.0mm}
					\item Inputs are points on the unit circle with uniform distribution:
					      \[
						      x(\varphi)=(\cos\varphi,\sin\varphi)\in S^1.
					      \]

					\item Fourier basis functions are written compactly as
					      \[
						      \phi_k \in \{1\}\cup\{\sin(k\varphi),\cos(k\varphi): k\ge1\}.
					      \]

					\item The target is a Fourier sum with amplitudes \(c_k\), the weights of each mode:
					      \[
						      f^\star(\varphi)=\sum_k c_k\phi_k(\varphi),
						      \qquad
						      r_t=f_t-f^\star=\sum_k\alpha_k(t)\phi_k.
					      \]

					\item The model is a one-hidden-layer ReLU network of width \(m\), initialized with
					      \(w_i\sim\mathcal N(0,I_2)\), \(a_i\sim\mathcal N(0,1)\), and \(b_i=0\).

					\item \textbf{Question:} what controls the decay of each \(\alpha_k(t)\)?
				\end{itemize}
			}

			\vspace{1mm}
			\centering
			\begin{tikzpicture}[x=1cm,y=1cm,scale=0.78]
				\draw[line] (0,0) -- (4.4,0);
				\draw[line] (0,-0.55) -- (0,0.82);
				\draw[tudcyan,very thick,domain=0:6.283,samples=100,smooth,variable=\t]
				plot ({4.4*\t/6.283},
				{0.24*(1+0.9*cos(deg(\t))+0.7*cos(deg(3*\t))-0.5*sin(deg(5*\t))-0.4*sin(deg(6*\t))-0.5)});
				\node[font=\tiny,text=muted,anchor=west] at (4.46,0) {$\varphi$};
				\node[font=\tiny,text=muted,anchor=south] at (2.2,0.82) {target \(f^\star\)};
			\end{tikzpicture}
		\end{column}

		\begin{column}{0.44\textwidth}
			\centering
			\begin{tikzpicture}[scale=0.78]
				\draw[line] (-1.2,0) -- (1.2,0);
				\draw[line] (0,-1.2) -- (0,1.2);
				\node[font=\tiny,text=muted,anchor=west] at (1.02,-0.12) {$x_1$};
				\node[font=\tiny,text=muted,anchor=south west] at (0.10,0.98) {$x_2$};
				\draw[line] (0,0) circle (1.15);
				\draw[tudcyan,very thick] (0,0) -- ({1.15*cos(35)},{1.15*sin(35)});
				\fill[tudcyan] ({1.15*cos(35)},{1.15*sin(35)}) circle (1.8pt);
				\node[font=\tiny,text=muted] at (1.45,0.75) {$x(\varphi)$};
				\node[font=\tiny,text=muted] at (0,-1.45) {$S^1$};
				\node[font=\tiny,text=muted] at (0.55,0.18) {$\varphi$};
			\end{tikzpicture}

			\vspace{2mm}
			% Self-contained TikZ snippet — single-hidden-layer ReLU MLP on S^1.
% No preamble/document: meant to be pulled into a slide with \input.
% Uses the deck's palette (text, muted, line, border, tudcyan, bg) and
% \familydefault (sans), so it inherits the house style automatically.
\begin{tikzpicture}[
		x=1mm, y=1mm,
		every node/.style={font=\tiny},
		neuron/.style={circle, draw=text, fill=bg, minimum size=2mm, line width=0.5pt},
		outnode/.style={circle, draw=tudcyan, fill=bg, minimum size=2mm, line width=0.8pt},
		conn/.style={draw=text, line width=0.3pt},
		lbl/.style={font=\scriptsize, text=muted}
	]
	% --- node positions (mm); kept compact: natural size ~51mm x ~25mm,
	% i.e. roughly 38% of slide width and ~30mm tall (incl. labels) as required ---
	\coordinate (i1) at (0, 5);
	\coordinate (i2) at (0, -5);
	\coordinate (h1) at (14, 10);
	\coordinate (h2) at (14, 4);
	\coordinate (h3) at (14, -4);
	\coordinate (h4) at (14, -10);
	\coordinate (o)  at (28, 0);

	% --- edges first, so the nodes sit cleanly on top of them ---
	\foreach \a in {i1, i2}
	\foreach \b in {h1, h2, h3, h4}
	\draw[conn] (\a) -- (\b);
	\foreach \b in {h1, h2, h3, h4}
	\draw[conn] (\b) -- (o);

	% --- nodes (drawn over the edges; fill=bg masks the stubs cleanly) ---
	\node[neuron] at (i1) {};
	\node[neuron] at (i2) {};
	\node[neuron] at (h1) {};
	\node[neuron] at (h2) {};
	\node[text=text, font=\small] at (14, 1) {$\vdots$};
	\node[neuron] at (h3) {};
	\node[neuron] at (h4) {};
	\node[outnode] at (o) {};

	% --- labels (absolute coords, so no \usetikzlibrary{calc} is needed and
	% the snippet stays self-contained against the deck's default tikz setup) ---
	\node[lbl, anchor=east] at (-2, 0) {$x(\varphi)$};
	\node[lbl, anchor=west, text=tudcyan] at (30.5, 0) {$f(x)$};
	\node[lbl, align=center] at (14, -14) {$m$ neurons};
\end{tikzpicture}

\par\vspace{1mm}
{\scriptsize\itshape\textcolor{muted}{%
		$\displaystyle f_m(x;\theta)=\frac{1}{\sqrt m}\sum_{i=1}^m a_i\,\sigma\!\left(w_i^\top x + b_i\right),\ \ \sigma=\mathrm{ReLU}$%
	}}

		\end{column}
	\end{columns}
\end{frame}

% ------------------------------------------------------------------
% Reference model: the ideal infinite-width / continuum-data operator.
% Message: convolution on the circle means Fourier modes diagonalize the
% dynamics; eigenvalues give learning rates.
\begin{frame}
	\note<1->{
		\begin{itemize}
			\item Infinite width and continuum data.
			\item Uniform circle makes \(T\) a convolution operator.
			\item Fourier modes diagonalize convolution.
			\item Each mode decays independently at rate \(\lambda_k\).
			\item Formula shows \(\lambda_k\sim k^{-2}\): high frequencies are slow.
		\end{itemize}
	}
	\vspace{0.2em}
	\slidetitle{REFERENCE MODEL}{The ideal picture: a convolution operator on $S^1$}
	{\textcolor{muted}{At infinite width and with continuum data, the NTK becomes fixed and rotation-invariant.}}

	{\raggedleft\tiny\textcolor{muted}{recall: fixed \(T\) gives \(\dot\alpha_k=-\lambda_k\alpha_k\)}\par}

	\vspace{0.5mm}
	\begin{whitecard}
		{\scriptsize
			With uniform measure \(\mu(d\psi)=d\psi/(2\pi)\), the fixed operator acts by
		}\par\vspace{-0.5mm}
		{\scriptsize
			\[
				(T r)(\varphi)=
				\int_{0}^{2\pi}\Theta(\varphi-\psi)\,r(\psi)\,\frac{d\psi}{2\pi}.
			\]
		}
		\vspace{-1mm}
		{\scriptsize
			Because this is a convolution on the circle, the Fourier modes are eigenfunctions:
		}\par
		{\scriptsize
		\[
			T\phi_k = \lambda_k \phi_k,
			\qquad
			\alpha_k(t)=e^{-\lambda_k t}\alpha_k(0).
		\]
		}
	\end{whitecard}

	\vspace{0.5mm}
	\begin{columns}[T,onlytextwidth]
		\begin{column}{0.53\textwidth}
			\begin{whitecard}[height=35mm]
				{\scriptsize\bfseries Explicit eigenvalues}\hfill{\tiny\textcolor{muted}{Thm.~5.2}}\par\vspace{0.5mm}
				{\scriptsize
					\begingroup
					\renewcommand{\arraystretch}{1.5}
					\[
						\lambda_k =
						\left\{
						\begin{array}{@{}ll@{}}
							\frac14+\frac{3}{\pi^2},
							 & k=0,                 \\[2mm]
							\frac14+\frac{1}{\pi^2},
							 & k=1,                 \\[2mm]
							\frac{k^2+3}{\pi^2(k^2-1)^2},
							 & k\ge 2 \text{ even}, \\[2mm]
							\frac{1}{\pi^2 k^2},
							 & k\ge 3 \text{ odd}.
						\end{array}
						\right.
					\]
					\endgroup
				}
			\end{whitecard}
		\end{column}

		\begin{column}{0.43\textwidth}
			\begin{bluecard}[height=17mm]
				{\scriptsize\bfseries Interpretation}\par\vspace{0.5mm}
				\begin{itemize}\tightlist\scriptsize
					\item Low frequencies $\implies $ Small \(k\): larger \(\lambda_k\), faster decay.
					\item High frequencies $\implies$ Large \(k\): \(\lambda_k\sim k^{-2}\), slower decay.
				\end{itemize}
			\end{bluecard}

			\vspace{1mm}
			\begin{whitecard}[height=11mm]
				\centering
				{\scriptsize\bfseries This exactly predicts the demo:}
				{\scriptsize coarse structure first, fine detail later.}
			\end{whitecard}
		\end{column}
	\end{columns}

	% \vspace{-0.5mm}
	{\tiny\textcolor{muted}{
		\href{https://arxiv.org/abs/1806.07572}{Jacot et al., 2018};
		\href{https://arxiv.org/abs/1912.01198}{Cao et al., 2019}
	}}
\end{frame}

% ------------------------------------------------------------------
\begin{frame}
	\note<1->{
		\begin{itemize}
			\item The ideal model is useful only if we understand what happens when its assumptions are relaxed.
			\item Perturbation 1: finite samples.
			\item Replace continuum measure by random sample points:
			      \[
				      \mu \rightsquigarrow \frac1n\sum_i\delta_{\varphi_i}.
			      \]
			\item This breaks exact Fourier orthogonality and exact diagonalization.
			\item Question: do Fourier modes still behave approximately like eigenvectors of the sampled operator?
			\item Perturbation 2: frozen finite width.
			\item Replace deterministic $T$ by random finite-width tangent operator $T_0^{(m)}$.
			\item Question: are Fourier eigenspaces and learning rates preserved?
			\item Main theoretical part: control both finite data and frozen finite width.
		\end{itemize}
	}
	\slidetitle{THESIS QUESTION}{Does this spectral picture survive outside the ideal limit?}

	{\textcolor{muted}{The ideal model is useful only if we understand what happens when its assumptions are relaxed.}}

	\vspace{3mm}
	\begin{columns}[T,onlytextwidth]
		% ----------------------------------------------------------
		% Finite samples
		% ----------------------------------------------------------
		\begin{column}{0.45\textwidth}
			\visible<2->{
				\begin{yellowcard}[height=53mm]
					{\textcolor{orange}{\scriptsize\bfseries Perturbation I}}\par\vspace{1mm}
					{\bfseries Finite samples}\par\vspace{3mm}

					\[
						\mu
						\quad\rightsquigarrow\quad
						\frac1n\sum_{i=1}^n\delta_{\varphi_i}
					\]

					\vspace{1mm}
					{\scriptsize Replace the continuous measure by $n$ random points on $S^1$.}\par\vspace{3.5mm}

					\visible<3->{
						{\textcolor{orange}{\scriptsize\bfseries Question}}\par
						{\scriptsize Do Fourier modes still behave like eigenvectors of the sampled operator?}
					}
				\end{yellowcard}
			}
		\end{column}

		% ----------------------------------------------------------
		% VS marker
		% ----------------------------------------------------------
		\begin{column}{0.06\textwidth}
			\visible<4->{
				\centering\vspace{28mm}
				\tikz{\node[circle, draw=line, fill=white, inner sep=2.8pt, font=\tiny\bfseries] {vs};}
			}
		\end{column}

		% ----------------------------------------------------------
		% Finite width
		% ----------------------------------------------------------
		\begin{column}{0.45\textwidth}
			\visible<4->{
				\begin{purplecard}[height=53mm]
					{\textcolor{purple}{\scriptsize\bfseries Perturbation II}}\par\vspace{1mm}
					{\bfseries Finite width}\par\vspace{3mm}

					\[
						T
						\quad\rightsquigarrow\quad
						T_0^{(m)}
					\]

					\vspace{1mm}
					{\scriptsize Replace the infinite-width operator by the finite-width tangent operator at initialization.}\par\vspace{3.5mm}

					\visible<5->{
						{\textcolor{purple}{\scriptsize\bfseries Question}}\par
						{\scriptsize Does the finite-width tangent operator preserve the Fourier eigenspaces and learning rates?}
					}
				\end{purplecard}
			}
		\end{column}
	\end{columns}

	\vspace{1.5mm}
	\visible<5->{
		{\centering\scriptsize\textcolor{muted}{Both perturbations break exact diagonalization; the goal is to quantify the finite-scale error (\textcolor{muted}{\ Bowman \& Mont\'ufar, 2022})}\par}
	}
\end{frame}

% ------------------------------------------------------------------
% Main results, message-first: on a fixed low-frequency Fourier window both
% finite-scale operators remain close to the ideal diagonal Fourier operator.
% Extra Gram/action/projection statements live in the backup frame.
\begin{frame}
	\note<1->{
		\begin{itemize}
			\item Fix one low-frequency Fourier subspace.
			\item On that subspace, the ideal operator is diagonal.
			\item Finite samples have an exponential tail in \(n\).
			\item Frozen finite width is unbiased and has an exponential tail in \(m\).
			\item Detailed Gram and action bounds are in backup.
		\end{itemize}
	}
	\slidetitle{NON-ASYMPTOTIC CONTROL}{Main results: the spectral picture survives finite scale}
	% {\textcolor{muted}{On one fixed low-frequency subspace, finite scale turns the ideal diagonal operator into a controlled perturbation.}}

	\vspace{0.5mm}
	\begin{whitecard}
		{\centering\scriptsize
			\(\mathcal H_K=\operatorname{span}\{1,\cos\varphi,\sin\varphi,\ldots,\cos(K\varphi),\sin(K\varphi)\}\),
			\quad \(d=2K+1\).\par\vspace{1mm}
			\(\Lambda_K=\operatorname{diag}(\lambda_0,\lambda_1,\lambda_1,\ldots,\lambda_K,\lambda_K)\),
			and \(T_K\) denotes the relevant finite-scale operator on \(\mathcal H_K\).\par}
	\end{whitecard}

	\vspace{1mm}
	\begin{columns}[T,onlytextwidth]
		\begin{column}{0.49\textwidth}
			\begin{yellowcard}[height=36mm]
				{\textcolor{orange}{\bfseries Finite samples}\hfill{\tiny\textcolor{muted}{Thm.~5.5}}}\par\vspace{1mm}
				{\scriptsize Let \(T_K\) be the sampled operator on \(\mathcal H_K\), using \(n\) uniform samples from \([0,2\pi)\). For every \(\varepsilon>0\),}\par\vspace{-1.5mm}
				{\scriptsize
					\[
						\mathbb P\!\left(
						\|T_K-\Lambda_K\|_{\max}\ge\varepsilon
						\right)
						\le
						2d^2\exp\!\left(-C n\varepsilon^2\right).
					\]
				}
				\vspace{-3mm}\par
				{\scriptsize Equivalently, with probability at least \(1-\delta\),}\par\vspace{-1.5mm}
				{\scriptsize
				\[
					\|T_K-\Lambda_K\|_{\max}
					=
					O\!\left(\sqrt{\frac{\log(d/\delta)}{n}}\right).
				\]
				}
			\end{yellowcard}
		\end{column}
		\pause
		\begin{column}{0.49\textwidth}
			\begin{purplecard}[height=36mm]
				{\textcolor{purple}{\bfseries Finite width}\hfill{\tiny\textcolor{muted}{Thm.~5.6}}}\par\vspace{1mm}
				{\scriptsize Let \(T_K^{(m)}\) be the width-\(m\) tangent operator on \(\mathcal H_K\). Then \(\mathbb E[T_K^{(m)}]=\Lambda_K\), and for every \(\varepsilon>0\),}\par\vspace{-1.5mm}
				{\scriptsize
				\[
					\mathbb P\!\left(
					\|T_K^{(m)}-\Lambda_K\|_{\max}\ge\varepsilon
					\right)
					\le
					2d^2\exp\!\left(-C m\varepsilon^2\right).
				\]
				}
				\vspace{-3mm}\par
				{\scriptsize Equivalently, with probability at least \(1-\delta\),}\par\vspace{-1.5mm}
				{\scriptsize
				\[
					\|T_K^{(m)}-\Lambda_K\|_{\max}
					=
					O\!\left(\sqrt{\frac{\log(d/\delta)}{m}}\right).
				\]
				}
			\end{purplecard}
		\end{column}
	\end{columns}

	\vspace{0.5mm}
	\pause
	\begin{whitecard}
		\centering
		{\scriptsize\bfseries What this gives us:}\ {\scriptsize a way to predict how the learning dynamics change with \(n\) and \(m\). In particular, the finite-sample block tells us how to precondition the operator and reduce the frequency-wise gap in learning rates.}
	\end{whitecard}
\end{frame}

% ------------------------------------------------------------------
\begin{frame}
	\note<1->{
		\begin{itemize}
			\item So far, the story has been: start from the ideal operator, then control finite-sample and frozen finite-width errors.
			\item But this still leaves out normal training.
			\item In normal training, the tangent kernel is allowed to move:
			      \[
				      T_0^{(m)} \rightsquigarrow T_t^{(m)}.
			      \]
			\item The question is whether that movement actually matters.
			\item At $m=128$ and $m=256$, it clearly does: the evolving kernel reaches lower training error than the frozen one.
			\item At $m=1024$, the reversal appears: the frozen kernel reaches the lower final loss.
			\item So kernel evolution seems to be a genuine small-width effect.
			\item This is the main tension: finite width can improve learning by changing the operator, but it also breaks the clean spectral reference model.
		\end{itemize}
	}
	\slidetitle{EVOLVING KERNEL}{A finite-width effect not explained by the NTK at initialization}
	{\textcolor{muted}{After controlling the operator at initialization, the next question is whether training changes the operator enough to matter.}}

	\pause
	\begin{columns}[T,onlytextwidth]
		\begin{column}{0.49\textwidth}
			\begin{whitecard}
				\centering
				\includegraphics[width=\linewidth]{train_ls_frozen_vs_evolving_w128.png}\par
				{\tiny $m=128$: evolving reaches lower training error}
			\end{whitecard}
		\end{column}
		\pause
		\begin{column}{0.49\textwidth}
			\begin{whitecard}
				\centering
				\includegraphics[width=\linewidth]{train_ls_frozen_vs_evolving_w1024.png}\par
				{\tiny $m=1024$: frozen reaches lower final training error}
			\end{whitecard}
		\end{column}
	\end{columns}
\end{frame}

% ------------------------------------------------------------------
\begin{frame}
	\note<1->{
		The first possible explanation is geometric: maybe training helps because the
		evolving kernel aligns its eigenspaces better with the Fourier modes. These
		plots rule out that simple story. At width 128, the low modes remain aligned,
		but the higher-frequency blocks lose alignment. At width 1024, the frozen
		initialization is already very well aligned, and the evolving kernel again
		moves several higher-frequency blocks away from the Fourier reference. So
		the small-width improvement is not coming from better Fourier alignment.
	}
	\slidetitle{EVOLVING KERNEL}{First check: does alignment improve?}
	{\textcolor{muted}{If evolution helped by preserving Fourier geometry, the alignment curves should rise above their initial alignment. They do not.}}

	\vspace{1mm}
	\begin{columns}[T,onlytextwidth]
		\begin{column}{0.495\textwidth}
			\begin{whitecard}[left=0.7mm,right=0.7mm,top=0.7mm,bottom=0.9mm]
				\centering
				\includegraphics[width=\linewidth,height=39mm,keepaspectratio]{rms_cos_alignment_by_freq_w128.png}\par
				{\tiny $m=128$: higher-frequency alignment drops sharply.}
			\end{whitecard}
		\end{column}
		\begin{column}{0.495\textwidth}
			\begin{whitecard}[left=0.7mm,right=0.7mm,top=0.7mm,bottom=0.9mm]
				\centering
				\includegraphics[width=\linewidth,height=39mm,keepaspectratio]{rms_cos_alignment_by_freq_w1024.png}\par
				{\tiny $m=1024$: evolution still worsens high frequencies.}
			\end{whitecard}
		\end{column}
	\end{columns}

	\vspace{1mm}
	\begin{bluecard}[height=8mm,valign=center]
		\centering
		{\scriptsize\bfseries Alignment does not explain why the evolving kernel reaches lower training error at lower widths.}
	\end{bluecard}
\end{frame}

% ------------------------------------------------------------------
\begin{frame}
	\note<1->{
		The second explanation is strength. Even if the eigenspaces do not become
		better aligned with the Fourier modes, the operator can still become stronger
		on some low-frequency blocks. This is what the spectral-strength plots suggest.
		At width 128, the constant block and first frequency increase sharply, with a
		smaller increase for k equals 2. At width 1024, the same pattern is visible
		but much weaker. This does not prove a mechanism, but it suggests that the
		low-frequency preference is still involved in the evolving-kernel behaviour.
	}
	\slidetitle{EVOLVING KERNEL}{What changes: low-frequency strength rises}
	{\textcolor{muted}{The moving kernel becomes stronger on some low-frequency blocks, even without better Fourier alignment.}}

	\vspace{1mm}
	\begin{columns}[T,onlytextwidth]
		\begin{column}{0.49\textwidth}
			\begin{whitecard}
				\centering
				\includegraphics[width=\linewidth,height=30mm,keepaspectratio]{spectral_strength_by_freq_w128.png}\par
				{\tiny $m=128$: large increase in $k=0$ and $k=1$.}
			\end{whitecard}
		\end{column}
		\begin{column}{0.49\textwidth}
			\begin{whitecard}
				\centering
				\includegraphics[width=\linewidth,height=30mm,keepaspectratio]{spectral_strength_by_freq_w1024.png}\par
				{\tiny $m=1024$: same direction of change, but much weaker.}
			\end{whitecard}
		\end{column}
	\end{columns}

	\vspace{1mm}
	\begin{bluecard}
		\centering
		{\bfseries The low-frequency preference is still visible, but the mechanism is not fully explained.}
	\end{bluecard}

	\vspace{0.7mm}
	\begin{whitecard}[height=8mm,valign=center]
		\centering
		{\tiny This is a diagnostic correlation, not a causal explanation.\quad
			The effect is strongest at small width and much weaker at $m=1024$.}
	\end{whitecard}
\end{frame}

% ------------------------------------------------------------------
\begin{frame}
	\note<1->{ Map each answer back to the three questions posed early in the talk. }
	\slidetitle{CONCLUSION}{What we can now say}
	{\textcolor{muted}{The three questions from the start, answered.}}

	\vspace{3mm}
	{\small
		\begin{itemize}
			\setlength{\itemsep}{4mm}
			\item[\questionnum{blue}{1}] {\bfseries What sets the learning speed of each frequency?}
			      \begin{itemize}\tightlist\scriptsize
				      \item Mode-wise rates are eigenvalues of the NTK; $\lambda_k\sim k^{-2}$, so low frequencies are learned first.
			      \end{itemize}

			\item[\questionnum{green}{2}] {\bfseries Does this survive with only $n$ sampled points?}\enspace{\scriptsize\textcolor{green}{\emph{RQ1}}}
			      \begin{itemize}\tightlist\scriptsize
				      \item Yes. Finite samples preserve the picture with explicit $O(n^{-1/2})$ control.
			      \end{itemize}

			\item[\questionnum{orange}{3}] {\bfseries Does it survive with finitely many neurons, at init and during training?}\enspace{\scriptsize\textcolor{orange}{\emph{RQ2}}}
			      \begin{itemize}\tightlist\scriptsize
				      \item At initialisation, yes, with explicit $O(m^{-1/2})$ control. Once the kernel evolves the clean spectral picture no longer fully describes training: we observe increased low-frequency spectral strength alongside lower training error, but the cause is still open.
			      \end{itemize}
		\end{itemize}
	}
\end{frame}

% ------------------------------------------------------------------
\begin{frame}
	\note<1->{
		The exact Fourier picture rests on four assumptions, and each one tells us
		what changes when it is relaxed. The uniform measure is what makes the kernel
		a convolution, so the Fourier modes are its eigenfunctions; a non-uniform
		density or Gaussian inputs change the basis to Hermite polynomials, but
		eigenvalue size still sets the rates, and the proofs would need redoing for
		unbounded support. ReLU fixes the polynomial decay of the spectrum; a smooth
		activation keeps the same eigenbasis but decays faster. Freezing the kernel
		is what gives the concentration bounds; letting it evolve is the regime we
		only study empirically. And the bounds hold on a fixed low-frequency block,
		not uniformly over all modes. The point is that the framework is robust in
		direction, but the guarantees are local to these assumptions.
	}
	\slidetitle{DISCUSSION \& LIMITATIONS}{What we assumed, and what changes if we relax it}

	\vspace{2mm}
	{\renewcommand{\arraystretch}{1.35}
		\begin{tabular}{@{}p{0.18\linewidth}p{0.33\linewidth}p{0.43\linewidth}@{}}
			\toprule
			{\scriptsize\bfseries Assumption} & {\scriptsize\bfseries What it gives us}                                                                                                                                                                                                                   & {\scriptsize\bfseries If we relax it} \\
			\midrule
			{\scriptsize\textcolor{blue}{\bfseries Uniform measure on $S^1$}}
			                                  & {\scriptsize The NTK is a convolution, so the Fourier modes are exactly its eigenfunctions.}
			                                  & {\scriptsize A non-uniform density or Gaussian inputs change the basis (e.g.\ Hermite polynomials), but eigenvalue size still sets the rates. The bounded-kernel proofs assume a compact domain, so they need rederiving for unbounded support.}                                                  \\
			{\scriptsize\textcolor{green}{\bfseries ReLU activation}}
			                                  & {\scriptsize A closed-form spectrum with polynomial decay, $\lambda_k = O(k^{-2})$.}
			                                  & {\scriptsize A smooth activation keeps the Fourier eigenbasis but decays faster, sharpening the bias. The fixed-kernel analysis carries over with the spectrum recomputed.}                                                                                                                       \\
			{\scriptsize\textcolor{orange}{\bfseries Frozen (lazy) kernel}}
			                                  & {\scriptsize Finite width is a controlled perturbation, concentrating at $O(m^{-1/2})$.}
			                                  & {\scriptsize When the kernel evolves the initial spectrum no longer governs the dynamics; at small width we observe increased low-frequency strength alongside lower training error, with the cause still open. This regime is only studied empirically.}                                         \\
			{\scriptsize\textcolor{purple}{\bfseries Fixed block $\mathcal H_K$}}
			                                  & {\scriptsize A controlled entrywise error on a chosen low-frequency block.}
			                                  & {\scriptsize Over all modes, or with $K$ growing in $n,m$, the entrywise bounds are not sharp where eigenvalues are small and close together; operator-norm tools would tighten them.}                                                                                                            \\
			\bottomrule
		\end{tabular}
	}
\end{frame}

% ------------------------------------------------------------------
\begin{frame}
	\note<1->{
		Where this sits. The NTK turns a non-convex training problem into a linear
		one that can be characterised exactly, and our networks live between two
		regimes: lazy training, where the tangent kernel barely moves, and feature
		learning, where it moves substantially. Four points on where the work sits.
		First, it makes the NTK spectral-bias mechanism, usually argued
		asymptotically or shown empirically, explicit on the circle, where every rate
		is a computable eigenvalue. Second, the infinite-width NTK is often treated as
		a purely asymptotic object, and we instead give non-asymptotic bounds on a
		low-frequency block. Third, the fixed kernel is a useful baseline while it
		stays lazy, but it is not a full account once features are learned, and this
		work marks where that reduction stops. Fourth, the mechanism is the one
		practitioners already modify: PINNs rescale this spectrum, and Fourier
		features and SIREN reshape it. A result on the uniform circle is still far
		from a guarantee about a deployed model.
	}
	\slidetitle{DISCUSSION \& LIMITATIONS}{Where this sits, and where it connects}

	\vspace{2mm}
	{\small
		\begin{itemize}
			\setlength{\itemsep}{2.6mm}
			\item {\bfseries\textcolor{blue}{An exact, concrete case study}}
			      \begin{itemize}\tightlist\scriptsize
				      \item The NTK spectral-bias mechanism, usually argued asymptotically or shown empirically, made explicit on $S^1$, where every rate is a computable eigenvalue (\textcolor{muted}{\href{https://arxiv.org/abs/1806.07572}{Jacot et al., 2018}}; \textcolor{muted}{\href{https://arxiv.org/abs/1806.08734}{Rahaman et al., 2019}}).
			      \end{itemize}

			\item {\bfseries\textcolor{green}{Finite scale, not only the limit}}
			      \begin{itemize}\tightlist\scriptsize
				      \item The infinite-width NTK is often treated as a purely asymptotic object; we add non-asymptotic bounds on a low-frequency block, alongside work probing finite-width corrections and the kernel picture's limits (\textcolor{muted}{\href{https://arxiv.org/abs/1909.05989}{Hanin \& Nica, 2019}}; \textcolor{muted}{\href{https://arxiv.org/abs/2206.10012}{Vyas et al., 2022}}).
			      \end{itemize}

			\item {\bfseries\textcolor{orange}{A baseline, not the whole story}}
			      \begin{itemize}\tightlist\scriptsize
				      \item The fixed kernel predicts the rates while it stays lazy, but is not a full account once features are learned; this work marks where that reduction stops (\textcolor{muted}{\href{https://arxiv.org/abs/1812.07956}{Chizat et al., 2019}}).
			      \end{itemize}

			\item {\bfseries\textcolor{purple}{Connects to practice}}
			      \begin{itemize}\tightlist\scriptsize
				      \item The same spectrum PINNs rescale to fix high-frequency failure (\textcolor{muted}{\href{https://arxiv.org/abs/2007.14527}{Wang et al., 2022}}), and that Fourier features and SIREN reshape for coordinate networks (\textcolor{muted}{\href{https://arxiv.org/abs/2006.10739}{Tancik et al., 2020}}; \textcolor{muted}{\href{https://arxiv.org/abs/2006.09661}{Sitzmann et al., 2020}}).
			      \end{itemize}
		\end{itemize}
	}

	\vspace{1.5mm}
	{\scriptsize\textcolor{muted}{Still, a result on the uniform circle is far from a guarantee about a deployed model.}}
\end{frame}

% ------------------------------------------------------------------
\begin{frame}
	\note<1->{ One-line takeaway, then questions. }
	\vspace{8mm}
	\centering
	\begin{bluecard}[width=0.8\linewidth]
		\centering
		{\large\bfseries Fourier structure survives finite scale, but kernel evolution changes the story.}\par\vspace{2mm}
		{\scriptsize In the lazy regimes it predicts the rates; once the kernel evolves we observe increased low-frequency spectral strength alongside lower training error, but the cause is still open.}
	\end{bluecard}

	\vspace{6mm}
	{\Large\bfseries Questions?}

	\vspace{4mm}
	{\scriptsize\textcolor{muted}{[acknowledgements placeholder]}}
\end{frame}

% ==================================================================
\appendix
\begin{frame}[noframenumbering]
	\centering\vspace{2cm}
	{\Large\bfseries Backup}
\end{frame}

% ------------------------------------------------------------------
\begin{frame}[noframenumbering]
	\slidetitle{}{Continuum prediction: spectral bias}
	\vspace{1mm}
	\centering
	\includegraphics[width=0.92\linewidth,height=56mm,keepaspectratio]{story_01_spectral_bias_two_panel_first200.png}
\end{frame}

% ------------------------------------------------------------------
\begin{frame}[noframenumbering]
	\slidetitle{}{NTK: from parameter updates to an operator on functions}
	{\textcolor{muted}{Gradient flow on the weights induces a linear flow on the residual, governed by an operator built from the network's own gradients.}}

	\vspace{1.5mm}
	\begin{columns}[T,onlytextwidth]
		\begin{column}{0.5\textwidth}
			\begin{whitecard}[height=44mm]
				{\scriptsize\bfseries Setup}\par\vspace{1mm}
				{\scriptsize Least-squares loss on $L^2(\mu)$, with residual $r_t=f_t-y$:}\par\vspace{-1.5mm}
				{\scriptsize
				\[
					\mathcal L(\theta)=\tfrac12\|f_m(\cdot;\theta)-y\|_{L^2(\mu)}^2 .
				\]
				}
				\vspace{-3mm}\par
				{\scriptsize The empirical NTK is the gradient inner product}\par\vspace{-1.5mm}
				{\scriptsize
					\[
						\Theta_t^{(m)}(\varphi,\psi)=\bigl\langle\nabla_\theta f_t(\varphi),\nabla_\theta f_t(\psi)\bigr\rangle,
					\]
				}
				\vspace{-3mm}\par
				{\scriptsize with operator $\;(T_t^{(m)}f)(\varphi)=\int \Theta_t^{(m)}(\varphi,\psi)f(\psi)\,d\mu(\psi)$.}
			\end{whitecard}
		\end{column}
		\begin{column}{0.5\textwidth}
			\begin{bluecard}[height=44mm]
				{\scriptsize\bfseries Residual flow}\par\vspace{1mm}
				{\scriptsize Gradient flow on $\theta$ and the chain rule give a linear flow on the residual:}\par\vspace{-1mm}
				{\scriptsize
				\[
					\partial_t r_t=-T_t^{(m)}r_t,
					\qquad
					\tfrac{d}{dt}\tfrac12\|r_t\|^2=-\langle r_t,T_t^{(m)}r_t\rangle .
				\]
				}
				\vspace{-2mm}\par
				{\scriptsize\bfseries Two regimes inside one equation:}
				\begin{itemize}\tightlist\scriptsize
					\item \textcolor{blue}{Frozen}: replace $T_t^{(m)}$ by $T_0^{(m)}$ (kernel regression on the initial features).
					\item \textcolor{orange}{Evolving}: keep the full time dependence of $T_t^{(m)}$.
				\end{itemize}
				\vspace{0.5mm}
				{\tiny\textcolor{muted}{As $m\to\infty$, $T_0^{(m)}\to T$, the deterministic continuum operator.}}
			\end{bluecard}
		\end{column}
	\end{columns}
\end{frame}

% ------------------------------------------------------------------
% Precise versions of the finite-scale results (moved out of the main deck):
% matrix definitions, max-norm bounds, and the unbiasedness identities.
\begin{frame}[noframenumbering]
	\slidetitle{}{Finite-scale bounds: precise statements}
	{\textcolor{muted}{$\mathcal H_K$: span of the Fourier modes up to frequency $K$, $d=2K+1$; $\Phi$: those modes evaluated at the $n$ samples; $\|\cdot\|_{\max}$: entrywise max.}}
	{\tiny\textcolor{muted}{Main-slide notation: \(T_K^{(n)}\) corresponds to \(H\), and \(T_K^{(m)}\) corresponds to \(B_m^{(K)}\).}}

	\vspace{1mm}
	\begin{columns}[T,onlytextwidth]
		\begin{column}{0.52\textwidth}
			\begin{yellowcard}
				{\textcolor{orange}{\bfseries Finite samples}\hfill{\tiny\textcolor{muted}{Thm.~5.5}}}\par\vspace{1mm}
				{\scriptsize Sampled kernel, Gram, and compressed operator:}\par\vspace{-1.5mm}
				{\scriptsize
					\[
						A_{ij}=\tfrac1n T(\theta_i,\theta_j),
						\qquad
						G=\tfrac1n\Phi^\top\Phi,
						\qquad
						H=\tfrac1n\Phi^\top A\Phi
					\]
				}
				\vspace{-3.5mm}\par
				{\scriptsize With probability at least $1-\delta$:}\par\vspace{-1.5mm}
				{\scriptsize
					\[
						\begin{aligned}
							\|G-I\|_{\max}
							 & \le
							C_1\sqrt{\tfrac{\log(d/\delta)}{n}}
							 &     & \text{(orthogonality)} \\[0.6mm]
							\|H-\Lambda^{(n)}\|_{\max}
							 & \le
							C_2\sqrt{\tfrac{\log(d/\delta)}{n}}
							 &     & \text{(eigenvalues)}   \\[0.6mm]
							\|A\Phi-\Phi\Lambda^{(n)}\|_{\max}
							 & \le
							C_3\sqrt{\tfrac{\log(nd/\delta)}{n}}
							 &     & \text{(kernel action)}
						\end{aligned}
					\]
				}
				\vspace{-2.5mm}\par
				{\tiny\textcolor{muted}{$\Lambda^{(n)}$: ideal eigenvalues with an $O(1/n)$ correction; $\mathbb{E}[H]=\Lambda^{(n)}$ (Prop.~5.4).}}
			\end{yellowcard}
		\end{column}
		\begin{column}{0.46\textwidth}
			\begin{purplecard}
				{\textcolor{purple}{\bfseries Frozen finite width}\hfill{\tiny\textcolor{muted}{Thm.~5.6}}}\par\vspace{1mm}
				{\scriptsize Compress the width-$m$ tangent operator at initialization to $\mathcal H_K$:}\par\vspace{-1.5mm}
				{\scriptsize
				\[
					B_m^{(K)}=P_K T_0^{(m)} P_K,
					\qquad
					\Lambda_K=P_K T P_K
				\]
				}
				\vspace{-3.5mm}\par
				{\scriptsize Unbiased at every width, and concentrating:}\par\vspace{-1.5mm}
				{\scriptsize
				\[
					\mathbb{E}\,B_m^{(K)}=\Lambda_K
				\]
				}
				\vspace{-4mm}\par
				{\scriptsize
				\[
					\|B_m^{(K)}-\Lambda_K\|_{\max}
					\le
					C_4\sqrt{\tfrac{\log(d/\delta)}{m}}
				\]
				}
				\vspace{-2.5mm}\par
				{\tiny\textcolor{muted}{with probability at least $1-\delta$; sub-exponential regime omitted.}}
			\end{purplecard}
		\end{column}
	\end{columns}
\end{frame}

% ------------------------------------------------------------------
\begin{frame}[noframenumbering]
	\slidetitle{}{Finite samples: observed mode decay tracks the eigenvalue prediction}
	{\textcolor{muted}{$n=4096$ uniform samples. Solid: observed normalised amplitude $|\alpha_p(t)|/|\alpha_p(0)|$. Dashed: diagonal finite-sample prediction from $\lambda_p^{(n)}$.}}

	\vspace{1.5mm}
	\begin{columns}[T,onlytextwidth]
		\begin{column}{0.49\textwidth}
			\begin{whitecard}[left=0.8mm,right=0.8mm,top=1mm,bottom=1mm]
				\centering
				\includegraphics[width=\linewidth,height=37mm,keepaspectratio]{story_02_random_n4096_components_first200.png}\par
				{\tiny First 200 steps: low frequencies decay first, matching the prediction.}
			\end{whitecard}
		\end{column}
		\begin{column}{0.49\textwidth}
			\begin{whitecard}[left=0.8mm,right=0.8mm,top=1mm,bottom=1mm]
				\centering
				\includegraphics[width=\linewidth,height=37mm,keepaspectratio]{story_02_random_n4096_components_all_steps_high_freq_only.png}\par
				{\tiny Longer horizon: the prediction holds until a component becomes small, then deviates.}
			\end{whitecard}
		\end{column}
	\end{columns}

	\vspace{1mm}
	\begin{whitecard}[height=8mm,valign=center]
		\centering
		{\scriptsize Sampling turns the exact continuum diagonal dynamics into an approximate low-frequency description, with deviations once a mode has nearly decayed.}
	\end{whitecard}
\end{frame}

% ------------------------------------------------------------------
\begin{frame}[noframenumbering]
	\slidetitle{}{Preconditioning: rescale each retained mode by its eigenvalue}
	{\textcolor{muted}{The fixed-kernel dynamics are linear with per-mode rate $\lambda_p$. Rescaling the retained Fourier block flattens the rates. Theory: divide by predicted $\lambda_p^{(n)}$. Empirical: invert the observed block $C=G^{-1}H$.}}

	\vspace{1.5mm}
	\begin{columns}[T,onlytextwidth]
		\begin{column}{0.49\textwidth}
			\begin{whitecard}[left=0.8mm,right=0.8mm,top=1mm,bottom=1mm]
				\centering
				\includegraphics[width=\linewidth,height=37mm,keepaspectratio]{all_mode_decay_three_panel_mean_std_n4096.png}\par
				{\tiny Both preconditioners flatten the retained-mode decay; empirical is most uniform.}
			\end{whitecard}
		\end{column}
		\begin{column}{0.49\textwidth}
			\begin{whitecard}[left=0.8mm,right=0.8mm,top=1mm,bottom=1mm]
				\centering
				\includegraphics[width=\linewidth,height=37mm,keepaspectratio]{probe_target_vs_prediction_three_panel_mean_std_n4096.png}\par
				{\tiny The retained components reach a much smaller error on the probe grid.}
			\end{whitecard}
		\end{column}
	\end{columns}
\end{frame}

% ------------------------------------------------------------------
\begin{frame}[noframenumbering]
	\slidetitle{}{Preconditioning: the two floors that limit it}
	{\textcolor{muted}{The preconditioner is built only on the retained block, and is solved in finite precision. Each limit shows up as a different plateau.}}

	\vspace{1.5mm}
	\begin{columns}[T,onlytextwidth]
		\begin{column}{0.49\textwidth}
			\begin{whitecard}[left=0.8mm,right=0.8mm,top=1mm,bottom=1mm]
				\centering
				\includegraphics[width=\linewidth,height=37mm,keepaspectratio]{out_of_block_all_mode_decay_three_panel_mean_std_n4096.png}\par
				{\tiny Out-of-block modes ($k=10,12$) stay near their start; in-block modes plateau higher from coupling.}
			\end{whitecard}
		\end{column}
		\begin{column}{0.49\textwidth}
			\begin{whitecard}[left=0.8mm,right=0.8mm,top=1mm,bottom=1mm]
				\centering
				\includegraphics[width=\linewidth,height=37mm,keepaspectratio]{mode_decay_fp32_vs_fp64_three_panel_mean_std_n4096.png}\par
				{\tiny fp32 (solid) vs fp64 (dashed): the preconditioned plateau drops by orders of magnitude, so that floor is numerical ($\sim\kappa u$).}
			\end{whitecard}
		\end{column}
	\end{columns}

	\vspace{1mm}
	\begin{whitecard}[height=8mm,valign=center]
		\centering
		{\scriptsize The block correction is explicit and cheap, but acts only inside the retained block and bottoms out at the conditioning-set precision floor.}
	\end{whitecard}
\end{frame}

% ------------------------------------------------------------------
\begin{frame}[noframenumbering]
	\slidetitle{}{Concentration bounds: the constants, and where the magnitude bound is loose}
	\vspace{1mm}
	\begin{columns}[T,onlytextwidth]
		\begin{column}{0.49\textwidth}
			\begin{whitecard}
				\centering
				\includegraphics[width=\linewidth,height=37mm,keepaspectratio]{theory_envelope_compare_action_err_max.png}\par
				{\tiny finite samples: action error decreases with $n$}
			\end{whitecard}
		\end{column}
		\begin{column}{0.49\textwidth}
			\begin{whitecard}
				\centering
				\includegraphics[width=\linewidth,height=37mm,keepaspectratio]{03_concentration_B_m.png}\par
				{\tiny frozen finite width: block error decreases with $m$}
			\end{whitecard}
		\end{column}
	\end{columns}
\end{frame}

% ------------------------------------------------------------------
\begin{frame}[noframenumbering]
	\slidetitle{}{Frozen versus evolving: cross-width view}
	\vspace{1mm}
	\centering
	\includegraphics[width=0.92\linewidth,height=56mm,keepaspectratio]{cross_width_residual_decay.png}
\end{frame}

% ------------------------------------------------------------------
\begin{frame}[noframenumbering]
	\slidetitle{}{Kernel movement is a small-width effect}
	\vspace{1mm}
	\centering
	\includegraphics[width=0.86\linewidth,height=56mm,keepaspectratio]{full_subspace_block_drift_fro_Ct_minus_C0_by_width_full.png}
\end{frame}

% ------------------------------------------------------------------
\begin{frame}[noframenumbering]
	\slidetitle{}{Full per-subspace evolving-kernel diagnostics}
	\vspace{1mm}
	\centering
	\includegraphics[width=0.88\linewidth,height=59mm,keepaspectratio]{per_subspace_story_grid_absolute_w128_w256_w512_w1024.png}
\end{frame}

% ------------------------------------------------------------------
\begin{frame}[noframenumbering]
	\slidetitle{}{Computing the eigenvalues $\lambda_k$}
	{\textcolor{muted}{The eigenvalues are the Fourier coefficients of the convolution kernel $\Theta(\delta)$.}}

	\vspace{1.5mm}
	\begin{columns}[T,onlytextwidth]
		\begin{column}{0.52\textwidth}
			\begin{whitecard}[height=46mm]
				{\scriptsize\bfseries From gradients to a closed-form kernel}\par\vspace{1mm}
				{\scriptsize At init ($a_i,w_i\sim\mathcal N(0,I)$, $b_i=0$), the width-$m$ NTK concentrates on}\par\vspace{-1.5mm}
				{\scriptsize
					\[
						\Theta(x,y)=K_0+(x^\top y+1)K_1,
					\]
				}
				\vspace{-3mm}\par
				{\scriptsize with the ReLU expectations (angle $\delta$, $\cos\delta=x^\top y$)}\par\vspace{-1.5mm}
				{\scriptsize
					\[
						K_1=\tfrac{\pi-\delta}{2\pi},
						\quad
						K_0=\tfrac{1}{2\pi}\bigl(\sin\delta+(\pi-\delta)\cos\delta\bigr),
					\]
				}
				\vspace{-3mm}\par
				{\scriptsize so that}\par\vspace{-2mm}
				{\scriptsize
					\[
						\Theta(\delta)=\tfrac{1}{2\pi}\bigl(\sin\delta+2(\pi-\delta)\cos\delta+(\pi-\delta)\bigr).
					\]
				}
			\end{whitecard}
		\end{column}
		\begin{column}{0.46\textwidth}
			\begin{bluecard}[height=46mm]
				{\scriptsize\bfseries Convolution $\Rightarrow$ Fourier diagonal}\par\vspace{1mm}
				{\scriptsize $\Theta$ depends only on $\delta=\varphi-\psi$, so $T$ is a convolution and}\par\vspace{-1.5mm}
				{\scriptsize
					\[
						\lambda_k=\frac{1}{\pi}\int_0^\pi \Theta(\delta)\cos(k\delta)\,d\delta .
					\]
				}
				\vspace{-2.5mm}\par
				{\scriptsize The parity integral $\int_0^\pi(\pi-\delta)\cos(m\delta)\,d\delta=\tfrac{1-(-1)^m}{m^2}$ splits even and odd $k$:}\par\vspace{0.5mm}
				{\scriptsize
					\[
						\lambda_0=\tfrac14+\tfrac{3}{\pi^2},\quad
						\lambda_1=\tfrac14+\tfrac{1}{\pi^2},
					\]
					\[
						\lambda_k=\begin{cases}\tfrac{k^2+3}{\pi^2(k^2-1)^2}, & k\ge2\text{ even},\\[1mm]\tfrac{1}{\pi^2 k^2}, & k\ge3\text{ odd}.\end{cases}
					\]
				}
			\end{bluecard}
		\end{column}
	\end{columns}
\end{frame}

% ------------------------------------------------------------------
\begin{frame}[noframenumbering]
	\slidetitle{}{The $k=2$ alignment anomaly}
	{\textcolor{muted}{$\lambda_2$ follows the generic even-$k$ formula normally. The anomaly is in the geometry: at finite width, $k=2$ is unusually well aligned with its Fourier subspace.}}

	\vspace{2mm}
	\begin{columns}[T,onlytextwidth]
		\begin{column}{0.5\textwidth}
			\begin{whitecard}[height=40mm]
				{\scriptsize\bfseries Observation}\par\vspace{1mm}
				\begin{itemize}\tightlist\scriptsize
					\item Projector error grows with frequency, as expected, with one exception: $k=2$ has the smallest error among the displayed frequencies.
					\item It is also the best aligned, sitting above its neighbours rather than on the decreasing trend.
				\end{itemize}
			\end{whitecard}
		\end{column}
		\begin{column}{0.5\textwidth}
			\begin{bluecard}[height=40mm]
				{\scriptsize\bfseries A tentative explanation}\par\vspace{1mm}
				\begin{itemize}\tightlist\scriptsize
					\item $k=0$ gets a direct eigenvalue from the bias; $k=1$ is tied to the coordinate embedding and is the only odd mode in the no-bias kernel.
					\item The finite-width perturbation may therefore act differently on $k=0,1$ than on higher blocks, leaving $k=2$ relatively favoured.
				\end{itemize}
				\vspace{0.5mm}
				{\tiny\textcolor{muted}{This is an empirical observation, not a theoretical claim; the mechanism is not investigated further.}}
			\end{bluecard}
		\end{column}
	\end{columns}
\end{frame}

% ------------------------------------------------------------------
\begin{frame}[noframenumbering]
	\slidetitle{}{Eigenvalue ordering: where it holds, and what it depends on}
	{\textcolor{muted}{Because the domain and measure are rotation-symmetric, the eigenvalue ordering is also a frequency ordering. The trainable bias is what keeps it strict.}}

	\vspace{1.5mm}
	\begin{columns}[T,onlytextwidth]
		\begin{column}{0.46\textwidth}
			\begin{whitecard}[height=42mm]
				{\scriptsize\bfseries The spectrum is strictly decreasing}\par\vspace{1mm}
				{\scriptsize
					\begin{tabular}{@{}ll@{}}
						$\lambda_0\approx0.554$  & $\lambda_1\approx0.351$  \\
						$\lambda_2\approx0.079$  & $\lambda_3\approx0.011$  \\
						$\lambda_4\approx0.0086$ & $\lambda_5\approx0.0041$ \\
					\end{tabular}
				}\par\vspace{1.5mm}
				{\scriptsize Two interleaved subsequences (even and odd $k$), both $O(k^{-2})$, with $\lambda_k>\lambda_{k+1}$ throughout.}
			\end{whitecard}
		\end{column}
		\begin{column}{0.52\textwidth}
			\begin{orangecard}[height=42mm]
				{\scriptsize\bfseries What the ordering depends on}\par\vspace{1mm}
				\begin{itemize}\tightlist\scriptsize
					\item The gap between $k=0,1$ and the rest comes from the constant and the coordinate embedding.
					\item The odd modes $k\ge3$ are nonzero \emph{only} because of the trainable hidden bias. Without it, $\lambda^{(\mathrm{nb})}_k=0$ for odd $k\ge3$, so those modes are null and never learned by the fixed kernel.
					\item Even $k\ge2$ are unchanged by the bias.
				\end{itemize}
				\vspace{0.5mm}
				{\tiny\textcolor{muted}{Consistent with bias-free two-layer ReLU models not expressing odd harmonics beyond $k=1$.}}
			\end{orangecard}
		\end{column}
	\end{columns}
\end{frame}

% ------------------------------------------------------------------
\begin{frame}[noframenumbering]
	\slidetitle{}{Why $k=0,1$ gain strength: the parameterisation hypothesis}
	{\textcolor{muted}{With inputs $x(\varphi)=(\cos\varphi,\sin\varphi)$, the pre-activation $w_j^\top x+b_j=\|w_j\|\cos(\varphi-\psi_j)+b_j$ already contains $k=0,1$; higher harmonics appear only once the ReLU acts. The network may strengthen the directions it represents most directly.}}

	\vspace{1.5mm}
	{\scriptsize\bfseries Control: equal target amplitudes across all modes.}\par
	\vspace{1mm}
	\begin{columns}[T,onlytextwidth]
		\begin{column}{0.49\textwidth}
			\begin{whitecard}[left=0.8mm,right=0.8mm,top=1mm,bottom=1mm]
				\centering
				\includegraphics[width=\linewidth,height=33mm,keepaspectratio]{equal-amp/spectral_strength_by_freq_w128.png}\par
				{\tiny $m=128$: strength still concentrates on $k=0,1$.}
			\end{whitecard}
		\end{column}
		\begin{column}{0.49\textwidth}
			\begin{whitecard}[left=0.8mm,right=0.8mm,top=1mm,bottom=1mm]
				\centering
				\includegraphics[width=\linewidth,height=33mm,keepaspectratio]{equal-amp/spectral_strength_by_freq_w1024.png}\par
				{\tiny $m=1024$: same direction, much weaker.}
			\end{whitecard}
		\end{column}
	\end{columns}

	\vspace{1mm}
	\begin{bluecard}[height=8mm,valign=center]
		\centering
		{\scriptsize Even with equal target mass, the gain favours $k=0,1$, so it is not just tracking where the target places its mass.}
	\end{bluecard}
\end{frame}

% ------------------------------------------------------------------
\begin{frame}[noframenumbering]
	\slidetitle{}{A sharper probe: remove the lowest frequencies from the target}
	{\textcolor{muted}{Target restricted to $k\ge 3$ (no $k=0,1,2$ component). This separates two accounts: parameterisation (the gain stays on $k=0,1$, the directions the input embedding and bias represent directly) versus amplifying the largest residual (the gain follows the lowest \emph{present} frequency).}}

	\vspace{1.5mm}
	\begin{columns}[T,onlytextwidth]
		\begin{column}{0.49\textwidth}
			\begin{whitecard}[left=0.8mm,right=0.8mm,top=1mm,bottom=1mm]
				\centering
				\includegraphics[width=\linewidth,height=34mm,keepaspectratio]{no-k01/spectral_strength_by_freq_allfreq_w128.png}\par
				{\tiny $m=128$: $k=0,1$ still gain the most strength, though absent from the target.}
			\end{whitecard}
		\end{column}
		\begin{column}{0.49\textwidth}
			\begin{whitecard}[left=0.8mm,right=0.8mm,top=1mm,bottom=1mm]
				\centering
				\includegraphics[width=\linewidth,height=34mm,keepaspectratio]{no-k01/train_mse_frozen_vs_evolving_w128.png}\par
				{\tiny $m=128$: evolving reaches lower training error than frozen, with $k=0,1,2$ removed.}
			\end{whitecard}
		\end{column}
	\end{columns}

	\vspace{1mm}
	\begin{bluecard}[height=11mm,valign=center]
		\centering
		{\tiny The strength gain still concentrates on $k=0,1$ although they are absent from the target, favouring the parameterisation account. This is one architecture on $S^1$ and a correlation, not an established cause.}
	\end{bluecard}
\end{frame}

% ------------------------------------------------------------------
\begin{frame}[noframenumbering]
	\slidetitle{}{Low-frequency strength gain is robust across amplitude configs}
	{\textcolor{muted}{Cumulative low-frequency strength $\mu_{\le k_{\max}}(t)$ for three amplitude profiles of the target: original, equal, and reversed (largest amplitudes on the highest frequencies).}}

	\vspace{1.5mm}
	\begin{columns}[T,onlytextwidth]
		\begin{column}{0.49\textwidth}
			\begin{whitecard}[left=0.8mm,right=0.8mm,top=1mm,bottom=1mm]
				\centering
				\includegraphics[width=\linewidth,height=37mm,keepaspectratio]{total_cumulative_mu_amp_configs_w128.png}\par
				{\tiny $m=128$: all three evolving configs rise well above the frozen baseline.}
			\end{whitecard}
		\end{column}
		\begin{column}{0.49\textwidth}
			\begin{whitecard}[left=0.8mm,right=0.8mm,top=1mm,bottom=1mm]
				\centering
				\includegraphics[width=\linewidth,height=37mm,keepaspectratio]{total_cumulative_mu_amp_configs_w1024.png}\par
				{\tiny $m=1024$: same ordering, smaller increase.}
			\end{whitecard}
		\end{column}
	\end{columns}

	\vspace{1mm}
	\begin{bluecard}[height=11mm,valign=center]
		\centering
		{\tiny The increase is largest for the equal-amplitude target and still clear for the reversed one, which puts more mass on high frequencies. So the low-frequency strength gain is not a consequence of the original target placing larger amplitudes on lower frequencies.}
	\end{bluecard}
\end{frame}

% ------------------------------------------------------------------
\begin{frame}[noframenumbering]
	\slidetitle{}{Mean-field viewpoint: a complementary large-width limit}
	{\textcolor{muted}{The NTK is one way to take width to infinity. The mean-field limit is another, and it is what makes the evolving-kernel regime natural.}}

	\vspace{2mm}
	\begin{columns}[T,onlytextwidth]
		\begin{column}{0.5\textwidth}
			\begin{bluecard}[height=42mm]
				{\scriptsize\textcolor{blue}{\bfseries NTK / lazy limit}}\par\vspace{1mm}
				\begin{itemize}\tightlist\scriptsize
					\item Track a \emph{fixed} tangent operator at initialization.
					\item On $S^1$ it is diagonal in the Fourier basis, so learning is read off eigenvalues and eigenspaces.
					\item Features barely move; this is the regime our bounds control.
				\end{itemize}
			\end{bluecard}
		\end{column}
		\begin{column}{0.5\textwidth}
			\begin{orangecard}[height=42mm]
				{\scriptsize\textcolor{orange}{\bfseries Mean-field limit}}\par\vspace{1mm}
				{\scriptsize With scaling $f_\theta=\tfrac1m\sum_\alpha v_\alpha\varphi(w_\alpha^\top x)$, the output depends on the neurons only through their empirical measure}\par\vspace{-1.5mm}
				{\scriptsize
					\[
						\mu_m=\tfrac1m\textstyle\sum_\alpha \delta_{\theta_\alpha}\;\longrightarrow\;\mu_t .
					\]
				}
				\vspace{-2mm}\par
				\begin{itemize}\tightlist\scriptsize
					\item Training is a flow of the \emph{distribution} of neurons, not a fixed kernel.
					\item This is the language for the feature-learning side where the kernel evolves.
				\end{itemize}
			\end{orangecard}
		\end{column}
	\end{columns}
\end{frame}

% ------------------------------------------------------------------
\begin{frame}[noframenumbering,t]
	\slidetitle{}{References}
	\vspace{2mm}
	{\scriptsize
		\begin{itemize}\tightlist
			\item Rahaman, Baratin, Arpit, Draxler, Lin, Hamprecht, Bengio, Courville (2019).
			      \emph{On the Spectral Bias of Neural Networks}. ICML.
			\item Jacot, Gabriel, Hongler (2018).
			      \emph{Neural Tangent Kernel: Convergence and Generalization in Neural Networks}. NeurIPS.
			\item Bowman, Mont\'ufar (2022).
			      \emph{Spectral Bias Outside the Training Set for Deep Networks in the Kernel Regime}. NeurIPS.
			\item Zhang, Bengio, Hardt, Recht, Vinyals (2017).
			      \emph{Understanding Deep Learning Requires Rethinking Generalization}. ICLR.
		\end{itemize}
	}
	\vfill
\end{frame}

\end{document}
